\chapter{Recursive Domain Equations by Finite-Flat Cofinal Approximation}
\label{ch:recursive-domain-equations-finite-flat}

In the preceding chapters we constructed the fuzzy ideal monad
\[
    T_{\mathbb V}X
    =
    \operatorname{Idl}_{\mathbb V}(X),
\]
where \(\operatorname{Idl}_{\mathbb V}(X)\) is the full
\(\mathbb V\)-subcategory of
\[
    [X^{\op},\underline{\mathbb V}]
\]
spanned by the finite-flat weights
\[
    I:X^{\op}\to\underline{\mathbb V}.
\]
Here \(\mathbb V\) denotes the fuzzy cosmos and
\(\underline{\mathbb V}\) denotes its self-enriched
\(\mathbb V\)-category of truth values.

Strict algebras for this monad are the fuzzy domains.  Equivalently, a fuzzy
domain is a small \(\mathbb V\)-category equipped with a strict coherent choice
of finite-flat weighted colimits, and a fuzzy-domain morphism is a
finite-flat-colimit-preserving \(\mathbb V\)-functor.

We also use the symmetric monoidal closed structure on fuzzy domains:
\[
    A\boxtimes_{\mathrm{ff}}B
    \dashv
    [B,-]_{\mathrm{ff}}.
\]
Thus
\[
    \mathbf{FDom}_{\mathbb V}
    (A\boxtimes_{\mathrm{ff}}B,C)
    \cong
    \mathbf{FDom}_{\mathbb V}
    \bigl(A,[B,C]_{\mathrm{ff}}\bigr),
\]
where \([B,C]_{\mathrm{ff}}\) is the fuzzy domain of
finite-flat-colimit-preserving \(\mathbb V\)-functors \(B\to C\).

The purpose of this chapter is to formulate recursive domain equation
constructions inside the finite-flat doctrine.  The classical construction by
\(\omega\)-chain bilimits in dcpos is recovered as the crisp separated
specialization.  The general replacement for an ordinary directed indexing
chain is a finite-flat approximation shape
\[
    (A,W),
    \qquad
    W:A^{\op}\to\underline{\mathbb V}.
\]
Cofinality is expressed by lower direct image of weights.  Tails and common
successors are not chosen arbitrarily; they are given by canonical
comma-style constructions, and the relevant exactness conditions say that
these canonical constructions are cofinal in the finite-flat sense.

The semantic target in the general resource-sensitive setting is not an
ordinary cartesian reflexive object
\[
    D\simeq [D,D]_{\mathrm{ff}}.
\]
In an affine enriched setting there is no canonical diagonal functor
\[
    A\to A\tensor A,
\]
since such a functor would require canonical contraction morphisms
\[
    A(a,a')\to A(a,a')\tensor A(a,a').
\]
Affinity gives weakening morphisms
\[
    A(a,a')\to I,
\]
not contraction into tensor squares.

Consequently the general internal-hom approximation theorem is
external-product-indexed:
\[
    (A\tensor A,W\boxtimes W).
\]
The same-index collapse is recovered only in the crisp separated/cartesian
specialization, where the ordinary diagonal
\[
    \omega\to\omega\times\omega
\]
exists and is cofinal for the ordinary directed ideal.

\section{Finite-flat approximation shapes}
\label{sec:finite-flat-approximation-shapes}

Throughout this chapter, finite-flatness is understood relative to the
finite-limit doctrine used in the definition of categorical fuzzy frames.  Thus
in a predicate category \([X,\underline{\mathbb V}]\), the specified finite
limits are computed pointwise.  In the posetal scalar-frame specialization this
means preservation of the terminal predicate and binary meets.

\begin{definition}[Finite-flat approximation shape]
\label{def:finite-flat-approximation-shape}
A \emph{finite-flat approximation shape} is a pair
\[
    (A,W),
\]
where \(A\) is a small \(\mathbb V\)-category and
\[
    W:A^{\op}\to\underline{\mathbb V}
\]
is a finite-flat weight.  Equivalently,
\[
    W\in\operatorname{Idl}_{\mathbb V}(A).
\]

If \(X\) is a fuzzy domain and
\[
    D:A\to X
\]
is a \(\mathbb V\)-functor, the corresponding approximation object is the
finite-flat weighted colimit
\[
    W*D.
\]
\end{definition}

Thus the passage from ordinary directed joins to fuzzy-domain approximation is
the passage
\[
    \bigvee_{a\in A}D_a
    \quad\leadsto\quad
    W*D.
\]

\begin{definition}[Lower direct image of a weight]
\label{def:lower-direct-image-weight-rde}
Let
\[
    h:A\to B
\]
be a \(\mathbb V\)-functor, and let
\[
    W:A^{\op}\to\underline{\mathbb V}
\]
be a weight.  The lower direct image of \(W\) along \(h\) is the weight
\[
    h_!W:B^{\op}\to\underline{\mathbb V}
\]
defined by
\[
    (h_!W)(b)
    =
    \int^{a\in A}
        W(a)\tensor B(b,ha).
\]
Equivalently, \(h_!W\) is the weighted left Kan extension of \(W\) along
\[
    h^{\op}:A^{\op}\to B^{\op}.
\]
\end{definition}

By the finite-flat direct-image theorem for the ideal doctrine, if \(W\) is
finite-flat, then \(h_!W\) is finite-flat.

\begin{definition}[Cofinal map of approximation shapes]
\label{def:cofinal-map-approximation-shapes}
Let
\[
    (A,W),
    \qquad
    (B,U)
\]
be finite-flat approximation shapes.  A \(\mathbb V\)-functor
\[
    h:A\to B
\]
is \emph{cofinal from \((A,W)\) to \((B,U)\)} if the canonical comparison
exhibits an isomorphism of weights
\[
    h_!W\cong U.
\]
When \(A=B\) and \(U=W\), we say that \(h\) is \(W\)-cofinal.
\end{definition}

Thus cofinality is equality of approximation weights after lower direct image,
not merely a pointwise finality condition on objects.

\begin{proposition}[Cofinality for finite-flat weighted colimits]
\label{prop:finite-flat-cofinality}
Let
\[
    h:(A,W)\to(B,U)
\]
be a cofinal map of finite-flat approximation shapes.  Then for every fuzzy
domain \(X\) and every diagram
\[
    D:B\to X,
\]
there is a canonical isomorphism
\[
    U*D
    \cong
    W*(Dh).
\]
\end{proposition}

\begin{proof}
Since \(h\) is cofinal,
\[
    U\cong h_!W.
\]
For every \(x\in X\), the universal property of weighted colimits and the
weighted left-Kan-extension adjunction give
\[
\begin{aligned}
    X((h_!W)*D,x)
    &\cong
    [B^{\op},\underline{\mathbb V}]
    (h_!W,X(D-,x))
    \\
    &\cong
    [A^{\op},\underline{\mathbb V}]
    (W,X(Dh-,x))
    \\
    &\cong
    X(W*(Dh),x).
\end{aligned}
\]
By enriched Yoneda,
\[
    (h_!W)*D\cong W*(Dh).
\]
Using \(U\cong h_!W\), this gives
\[
    U*D\cong W*(Dh).
\]
\end{proof}

\begin{definition}[External product of approximation shapes]
\label{def:external-product-approximation-shapes-rde}
Let
\[
    (A,W),
    \qquad
    (B,Q)
\]
be finite-flat approximation shapes.  Their external product is
\[
    (A,W)\boxtimes(B,Q)
    =
    (A\tensor B,W\boxtimes Q),
\]
where
\[
    (W\boxtimes Q)(a,b)
    =
    W(a)\tensor Q(b).
\]
By closure of finite-flat weights under external products,
\[
    W\boxtimes Q
\]
is finite-flat.
\end{definition}

\begin{lemma}[Constant diagrams]
\label{lem:constant-diagrams-finite-flat}
Let
\[
    (A,W)
\]
be a finite-flat approximation shape.  For every fuzzy domain \(X\) and every
object \(x\in X\), the \(W\)-weighted colimit of the constant diagram at \(x\)
is canonically \(x\):
\[
    W*\Delta x\cong x.
\]
\end{lemma}

\begin{proof}
Let
\[
    !:A\to\mathbb 1
\]
be the terminal functor, and let
\[
    \bar x:\mathbb 1\to X
\]
classify the object \(x\).  The constant diagram at \(x\) is
\[
    \Delta x=\bar x!.
\]

Since \(W\) is finite-flat, its character
\[
    \chi_W:[A,\underline{\mathbb V}]\to\underline{\mathbb V}
\]
preserves the terminal predicate.  Let
\[
    \top_A:A\to\underline{\mathbb V}
\]
be the constant upper predicate at \(I\).  Then
\[
\begin{aligned}
    (!_!W)(*)
    &=
    \int^{a\in A}W(a)\tensor \mathbb 1(*,!a)
    \\
    &\cong
    \int^{a\in A}W(a)\tensor I
    \\
    &=
    \chi_W(\top_A)
    \\
    &\cong
    I.
\end{aligned}
\]
Thus
\[
    !_!W\cong \mathbb 1(-,*).
\]
By Proposition~\ref{prop:finite-flat-cofinality},
\[
    W*(\bar x!)
    \cong
    (!_!W)*\bar x
    \cong
    \mathbb 1(-,*)*\bar x
    \cong
    x.
\]
\end{proof}

\begin{lemma}[Functoriality of finite-flat colimits in diagrams]
\label{lem:ff-colimits-functorial-in-diagrams}
Let \(C\) be a fuzzy domain, let
\[
    R:P^{\op}\to\underline{\mathbb V}
\]
be finite-flat, and let
\[
    H,H':P\to C
\]
be \(\mathbb V\)-functors.  A \(\mathbb V\)-natural transformation
\[
    \gamma:H\Rightarrow H'
\]
induces a canonical morphism
\[
    R*H\longrightarrow R*H'
\]
in \(C\).  This construction is functorial in \(\gamma\).
\end{lemma}

\begin{proof}
For each object \(c\in C\), the universal property of the weighted colimit gives
natural isomorphisms
\[
    C(R*H,c)
    \cong
    [P^{\op},\underline{\mathbb V}]
    (R,C(H-,c))
\]
and
\[
    C(R*H',c)
    \cong
    [P^{\op},\underline{\mathbb V}]
    (R,C(H'-,c)).
\]
The natural transformation
\[
    \gamma:H\Rightarrow H'
\]
induces by precomposition a \(\mathbb V\)-natural transformation
\[
    C(H'-,c)\Rightarrow C(H-,c).
\]
Transporting across the representing isomorphisms gives a morphism
\[
    C(R*H',c)\longrightarrow C(R*H,c)
\]
natural in \(c\).  By enriched Yoneda this family is represented by a unique
morphism
\[
    R*H\longrightarrow R*H'.
\]
Functoriality follows from functoriality of precomposition and uniqueness in
the representing property.
\end{proof}

\section{Support categories and projection-pair systems}
\label{sec:support-categories-split-systems}

Projection-pair systems indexed by an enriched category use global support
arrows for the actual transition projection pairs.  The remaining enriched
hom-structure enters through transition-strength morphisms.

\begin{definition}[Support category of an enriched index]
\label{def:support-category-enriched-index}
Let \(A\) be a small \(\mathbb V\)-category.  Its support category
\[
    A_0
\]
is the ordinary category with the same objects as \(A\), and hom-set
\[
    A_0(a,a')=\mathbb V(I,A(a,a')).
\]
If
\[
    u:I\to A(a,a')
    \qquad\text{and}\qquad
    v:I\to A(a',a'')
\]
are arrows in \(A_0\), their composite is the global element
\[
    I\cong I\tensor I
    \xrightarrow{\,v\tensor u\,}
    A(a',a'')\tensor A(a,a')
    \longrightarrow
    A(a,a'')
\]
induced by composition in \(A\).  Identities are the enriched identity maps
\[
    I\to A(a,a).
\]

In the posetal specialization, \(A_0(a,a')\) has a unique arrow precisely when
\[
    A(a,a')=\top.
\]
\end{definition}

\begin{lemma}[Support projections from an enriched tensor product]
\label{lem:support-projections-tensor-product}
Let \(A\) and \(B\) be small \(\mathbb V\)-categories.  There is a canonical
ordinary functor
\[
    (A\tensor B)_0\to A_0\times B_0.
\]
On objects it sends
\[
    (a,b)\longmapsto(a,b).
\]
On arrows, a global element
\[
    \xi:I\to A(a,a')\tensor B(b,b')
\]
is sent to the pair
\[
    \xi_A:I\to A(a,a'),
    \qquad
    \xi_B:I\to B(b,b'),
\]
where \(\xi_A\) is obtained by composing \(\xi\) with
\[
    A(a,a')\tensor B(b,b')
    \xrightarrow{1\tensor !}
    A(a,a')\tensor I
    \cong
    A(a,a'),
\]
and \(\xi_B\) is obtained similarly using
\[
    A(a,a')\tensor B(b,b')
    \xrightarrow{!\tensor 1}
    I\tensor B(b,b')
    \cong
    B(b,b').
\]
\end{lemma}

\begin{proof}
The terminality of \(I\) gives the weakening maps
\[
    B(b,b')\to I,
    \qquad
    A(a,a')\to I.
\]
Thus the two displayed projections are canonical.

Compatibility with identities is immediate from the unit constraints.  For
composition, the composition in \(A\tensor B\) is induced by the tensor of the
compositions in \(A\) and \(B\).  Applying the two weakening projections sends
this composite to the ordinary composites in \(A_0\) and \(B_0\).  Hence the
assignment defines an ordinary functor.
\end{proof}

\begin{definition}[The locally ordinary \(2\)-category of fuzzy domains]
\label{def:locally-ordinary-2-category-fdom}
Let
\[
    \mathbf{FDom}^{(2)}_{\mathbb V}
\]
be the locally ordinary \(2\)-category defined as follows.

Its objects are fuzzy domains.  Its \(1\)-cells are fuzzy-domain morphisms.
For fuzzy domains \(X,Y\), the hom-category is the support category of the
internal-hom fuzzy domain:
\[
    \mathbf{FDom}^{(2)}_{\mathbb V}(X,Y)
    =
    \bigl([X,Y]_{\mathrm{ff}}\bigr)_0.
\]
Thus a \(2\)-cell
\[
    \alpha:f\Rightarrow g
\]
is a global element
\[
    I\to [X,Y]_{\mathrm{ff}}(f,g).
\]

Horizontal composition is induced by the finite-flat-continuous composition
bimorphism
\[
    [Y,Z]_{\mathrm{ff}}\boxtimes_{\mathrm{ff}}[X,Y]_{\mathrm{ff}}
    \longrightarrow
    [X,Z]_{\mathrm{ff}}.
\]
Vertical composition is inherited from the support categories of the internal
homs.
\end{definition}

\begin{definition}[Fuzzy projection pair]
\label{def:fuzzy-projection-pair-ff}
Let \(X\) and \(Y\) be fuzzy domains.  A \emph{fuzzy projection pair}
\[
    (e,p):X\triangleleft Y
\]
consists of fuzzy-domain morphisms
\[
    e:X\to Y,
    \qquad
    p:Y\to X,
\]
together with an adjunction
\[
    e\dashv p
\]
in \(\mathbf{FDom}^{(2)}_{\mathbb V}\), whose unit
\[
    \eta:\id_X\Rightarrow pe
\]
is invertible.

Equivalently, after choosing the inverse of the unit, one has
\[
    pe\cong \id_X
\]
and a counit comparison
\[
    ep\Rightarrow \id_Y
\]
satisfying the triangle identities.  The morphism \(e\) is called the
embedding and \(p\) the projection.
\end{definition}

\begin{lemma}[Composition of fuzzy projection pairs]
\label{lem:composition-fuzzy-projection-pairs-ff}
Fuzzy projection pairs compose.  If
\[
    (e,p):X\triangleleft Y
\]
and
\[
    (e',p'):Y\triangleleft Z
\]
are fuzzy projection pairs, then
\[
    (e'e,pp'):X\triangleleft Z
\]
is a fuzzy projection pair.
\end{lemma}

\begin{proof}
The composite of adjunctions
\[
    e\dashv p,
    \qquad
    e'\dashv p'
\]
is an adjunction
\[
    e'e\dashv pp'.
\]
Its unit is the composite
\[
    \id_X
    \Rightarrow
    pe
    \Rightarrow
    pp'e'e.
\]
Both displayed unit maps are invertible, hence the composite unit is invertible.
The counit and the triangle identities are those of the composite adjunction.
\end{proof}

\begin{definition}[Projection-pair 2-cells]
\label{def:projection-pair-2-cells}
Let
\[
    (e,p),(e',p'):X\triangleleft Y
\]
be fuzzy projection pairs.  A 2-cell
\[
    (\alpha,\beta):(e,p)\Rightarrow(e',p')
\]
consists of enriched natural transformations
\[
    \alpha:e\Rightarrow e',
    \qquad
    \beta:p'\Rightarrow p
\]
which are mates under the adjunctions
\[
    e\dashv p,
    \qquad
    e'\dashv p'.
\]
\end{definition}

\begin{definition}[Projection-pair 2-category]
\label{def:projection-pair-2-category-ff}
The locally ordinary \(2\)-category
\[
    \mathbf{FDom}_{\mathbb V}^{\triangleleft}
\]
has fuzzy domains as objects, fuzzy projection pairs as \(1\)-cells, and
mate-compatible pairs of enriched natural transformations as \(2\)-cells.
The identity \(1\)-cell on \(X\) is
\[
    (\id_X,\id_X):X\triangleleft X,
\]
and composition is
\[
    (e',p')\circ(e,p)=(e'e,pp').
\]
Vertical and horizontal composition of \(2\)-cells are inherited from
\(\mathbf{FDom}^{(2)}_{\mathbb V}\); mate-compatibility is preserved by the
calculus of mates.
\end{definition}

\begin{definition}[Split \(A\)-indexed projection-pair system]
\label{def:split-A-indexed-projection-pair-system}
Let \(A\) be a small \(\mathbb V\)-category.  A
\emph{split \(A\)-indexed projection-pair system}
\[
    D:A\rightsquigarrow\mathbf{FDom}_{\mathbb V}^{\triangleleft}
\]
consists of:

\begin{enumerate}
    \item for every object \(a\in A\), a fuzzy domain \(D_a\);

    \item for every arrow
    \[
        u:a\to a'
    \]
    in the support category \(A_0\), a fuzzy projection pair
    \[
        (e_u,p_u):D_a\triangleleft D_{a'};
    \]

    \item normality isomorphisms
    \[
        (e_{1_a},p_{1_a})\cong(\id_{D_a},\id_{D_a});
    \]

    \item composition isomorphisms
    \[
        e_v e_u\cong e_{vu},
        \qquad
        p_u p_v\cong p_{vu}
    \]
    for composable support arrows
    \[
        a\xrightarrow{u}a'\xrightarrow{v}a'';
    \]

    \item the coherence equations for identities, composition, and
    associativity.
\end{enumerate}
\end{definition}

If
\[
    h:A\to B
\]
is a \(\mathbb V\)-functor, it induces an ordinary functor
\[
    h_0:A_0\to B_0.
\]
Thus a split \(B\)-indexed system
\[
    D:B\rightsquigarrow\mathbf{FDom}_{\mathbb V}^{\triangleleft}
\]
has a reindexing
\[
    Dh:A\rightsquigarrow\mathbf{FDom}_{\mathbb V}^{\triangleleft}
\]
defined on objects by
\[
    (Dh)_a=D_{ha}
\]
and on support arrows by applying \(D\) to their images under \(h_0\).

\begin{definition}[Compatible projection-pair cone]
\label{def:compatible-projection-pair-cone-split}
Let
\[
    D:A\rightsquigarrow\mathbf{FDom}_{\mathbb V}^{\triangleleft}
\]
be a split \(A\)-indexed projection-pair system, and let \(Z\) be a fuzzy
domain.

A \emph{compatible projection-pair cone} from \(D\) to \(Z\) consists of
projection pairs
\[
    (\varepsilon_a,\pi_a):D_a\triangleleft Z
    \qquad(a\in A),
\]
together with:

\begin{enumerate}
    \item for every support arrow
    \[
        u:a\to a'
    \]
    in \(A_0\), coherent isomorphisms
    \[
        \varepsilon_{a'}e_u\cong \varepsilon_a,
        \qquad
        p_u\pi_{a'}\cong \pi_a;
    \]

    \item writing
    \[
        r_a=\varepsilon_a\pi_a:Z\to Z,
    \]
    enriched transition-strength morphisms
    \[
        \theta^Z_{a,a'}:
        A(a,a')
        \longrightarrow
        [Z,Z]_{\mathrm{ff}}(r_a,r_{a'})
    \]
    satisfying the identity and composition axioms for a \(\mathbb V\)-functor
    \[
        r:A\to[Z,Z]_{\mathrm{ff}};
    \]

    \item for every support arrow
    \[
        u:I\to A(a,a'),
    \]
    the composite
    \[
        I
        \xrightarrow{u}
        A(a,a')
        \xrightarrow{\theta^Z_{a,a'}}
        [Z,Z]_{\mathrm{ff}}(r_a,r_{a'})
    \]
    is the canonical comparison
    \[
        \varepsilon_a\pi_a
        \cong
        \varepsilon_{a'}e_u p_u\pi_{a'}
        \Rightarrow
        \varepsilon_{a'}\pi_{a'}
    \]
    induced by the counit
    \[
        e_u p_u\Rightarrow \id_{D_{a'}}.
    \]
\end{enumerate}

A morphism of compatible cones is a family of mate-compatible enriched natural
transformations between the corresponding projection pairs, compatible with the
support-arrow cone isomorphisms and with the transition-strength morphisms.  The
resulting category of compatible cones is denoted
\[
    \operatorname{Cone}_{\triangleleft}(D;Z).
\]
\end{definition}

\begin{lemma}[The approximant functor of a compatible cone]
\label{lem:approximants-are-functorial}
Let
\[
    D:A\rightsquigarrow\mathbf{FDom}_{\mathbb V}^{\triangleleft}
\]
be a split \(A\)-indexed projection-pair system, and let
\[
    (\varepsilon_a,\pi_a):D_a\triangleleft Z
\]
be a compatible projection-pair cone.  Define
\[
    r_a=\varepsilon_a\pi_a:Z\to Z.
\]
Then the objects \(r_a\) assemble canonically to a \(\mathbb V\)-functor
\[
    r:A\to[Z,Z]_{\mathrm{ff}}.
\]
\end{lemma}

\begin{proof}
This is precisely the content of the transition-strength data in
Definition~\ref{def:compatible-projection-pair-cone-split}.
\end{proof}

\section{Canonical tails and common successors}
\label{sec:canonical-tails-common-successors}

The ordinary \(\omega\)-chain construction uses two elementary cofinality
facts: every tail of \(\omega\) is cofinal, and any two indices have a common
later index.  In the finite-flat enriched setting these facts are represented
by canonical support-comma constructions.

\subsection{Canonical tails}

\begin{definition}[Canonical support tail]
\label{def:canonical-support-tail}
Let \(A\) be a small \(\mathbb V\)-category and let \(a_0\in A\).  The
\emph{canonical support tail} of \(A\) at \(a_0\) is the enriched category
\[
    \operatorname{Tail}_A(a_0)
\]
defined as follows.

Its objects are pairs
\[
    (a,u),
    \qquad
    u:a_0\to a
\]
where \(u\) is an arrow in the support category \(A_0\).

The hom-object from
\[
    (a,u)
\]
to
\[
    (a',u')
\]
is the equalizer in \(\mathbb V\)
\[
    \operatorname{Tail}_A(a_0)((a,u),(a',u'))
    \longrightarrow
    A(a,a')
    \rightrightarrows
    A(a_0,a').
\]
The first parallel map is composition with the global element
\[
    u:I\to A(a_0,a).
\]
The second is the constant map determined by the global element
\[
    u':I\to A(a_0,a'),
\]
using the terminality of \(I\).  Composition and identities are inherited from
\(A\), by the universal property of the equalizers.

The evident projection
\[
    \tau_{a_0}:\operatorname{Tail}_A(a_0)\to A,
    \qquad
    (a,u)\longmapsto a,
\]
is a \(\mathbb V\)-functor.
\end{definition}

In the crisp preorder case, \(\operatorname{Tail}_A(a_0)\) is the full tail
suborder on the elements \(a\) satisfying \(a_0\leq a\).

\begin{definition}[Canonical tail weight]
\label{def:canonical-tail-weight}
Let
\[
    (A,W)
\]
be a finite-flat approximation shape and let \(a_0\in A\).  The
\emph{canonical tail weight} at \(a_0\) is
\[
    W^{a_0}
    =
    \tau_{a_0}^{*}W
    :
    \operatorname{Tail}_A(a_0)^{\op}
    \to
    \underline{\mathbb V}.
\]
Thus
\[
    W^{a_0}(a,u)=W(a).
\]
\end{definition}

\begin{definition}[Tail-cofinal approximation shape]
\label{def:tail-cofinal-approximation-shape}
A finite-flat approximation shape
\[
    (A,W)
\]
is \emph{tail-cofinal} if, for every object \(a_0\in A\), the canonical tail
weight
\[
    W^{a_0}=\tau_{a_0}^{*}W
\]
is finite-flat and the canonical direct-image comparison is an isomorphism:
\[
    (\tau_{a_0})_!W^{a_0}\cong W.
\]
\end{definition}

\begin{lemma}[Constant diagrams over canonical tails]
\label{lem:constant-diagrams-canonical-tail}
Let
\[
    (A,W)
\]
be tail-cofinal, and let \(a_0\in A\).  For every fuzzy domain \(X\) and every
object \(x\in X\),
\[
    W^{a_0}*\Delta x\cong x.
\]
\end{lemma}

\begin{proof}
By tail-cofinality, \(W^{a_0}\) is finite-flat.  Hence the result is
Lemma~\ref{lem:constant-diagrams-finite-flat} applied to
\[
    (\operatorname{Tail}_A(a_0),W^{a_0}).
\]
\end{proof}

\subsection{Canonical common successors}

\begin{definition}[Canonical common-successor category]
\label{def:canonical-common-successor-category}
Let \(P\) be a small \(\mathbb V\)-category.  The
\emph{canonical common-successor category}
\[
    \operatorname{Com}(P)
\]
is defined as follows.

An object of \(\operatorname{Com}(P)\) is a tuple
\[
    (i,j,k,u,v),
\]
where \(i,j,k\in P\), and
\[
    u:i\to k,
    \qquad
    v:j\to k
\]
are arrows in the support category \(P_0\).

For two objects
\[
    c=(i,j,k,u,v),
    \qquad
    c'=(i',j',k',u',v'),
\]
the hom-object
\[
    \operatorname{Com}(P)(c,c')
\]
is the equalizer in \(\mathbb V\) of the compatibility maps expressing
\[
    u'r=tu,
    \qquad
    v's=tv
\]
for triples
\[
    r:i\to i',
    \qquad
    s:j\to j',
    \qquad
    t:k\to k'.
\]
Explicitly, it is the equalizer of the two pairs of maps from
\[
    P(i,i')\tensor P(j,j')\tensor P(k,k')
\]
to
\[
    P(i,k')\times P(j,k')
\]
obtained respectively from the composites
\[
    u'r,\quad tu,
    \qquad
    v's,\quad tv.
\]
Composition and identities are inherited from those of \(P\).

There are canonical \(\mathbb V\)-functors
\[
    c_P:\operatorname{Com}(P)\to P\tensor P,
    \qquad
    m_P:\operatorname{Com}(P)\to P
\]
defined on objects by
\[
    c_P(i,j,k,u,v)=(i,j),
    \qquad
    m_P(i,j,k,u,v)=k.
\]
\end{definition}

In the crisp preorder case, an object of \(\operatorname{Com}(P)\) is exactly
a common upper bound
\[
    i\leq k,
    \qquad
    j\leq k.
\]

\begin{definition}[Canonical common-successor weight]
\label{def:canonical-common-successor-weight}
Let
\[
    (P,R)
\]
be a finite-flat approximation shape.  The
\emph{canonical common-successor weight} is
\[
    R^{\mathrm{com}}
    =
    m_P^{*}R:
    \operatorname{Com}(P)^{\op}
    \to
    \underline{\mathbb V}.
\]
Thus, on objects,
\[
    R^{\mathrm{com}}(i,j,k,u,v)=R(k).
\]
\end{definition}

\begin{definition}[Common-successor-cofinal approximation shape]
\label{def:common-successor-cofinal-shape}
A finite-flat approximation shape
\[
    (P,R)
\]
is \emph{common-successor-cofinal} if the canonical common-successor weight
\[
    R^{\mathrm{com}}=m_P^*R
\]
is finite-flat and the canonical maps
\[
    c_P:\operatorname{Com}(P)\to P\tensor P,
    \qquad
    m_P:\operatorname{Com}(P)\to P
\]
satisfy
\[
    (c_P)_!R^{\mathrm{com}}\cong R\boxtimes R,
    \qquad
    (m_P)_!R^{\mathrm{com}}\cong R.
\]
\end{definition}

Thus common-successor-cofinality is the finite-flat analogue of directedness:
double approximation indexed by \(P\tensor P\) can be compared with
approximation over canonical common successors, and then with single
approximation over \(P\).

\begin{definition}[Linear-cofinal finite-flat approximation shape]
\label{def:linear-cofinal-ff-approximation-shape}
A finite-flat approximation shape
\[
    (A,W)
\]
is \emph{linear-cofinal} if:

\begin{enumerate}
    \item it is tail-cofinal;

    \item its external product
    \[
        (A\tensor A,W\boxtimes W)
    \]
    is common-successor-cofinal.
\end{enumerate}
\end{definition}

This deliberately does not include a diagonal comparison
\[
    A\to A\tensor A.
\]
The diagonal is a cartesian feature and appears only in the crisp separated
collapse.

\section{Weight-cofinal iteration shapes}
\label{sec:weight-cofinal-iteration-shapes}

\begin{definition}[Weight-cofinal iteration shape]
\label{def:weight-cofinal-iteration-shape}
A \emph{weight-cofinal iteration shape} is a triple
\[
    (A,W,s),
\]
where \((A,W)\) is a finite-flat approximation shape and
\[
    s:A\to A
\]
is a \(W\)-cofinal endofunctor:
\[
    s_!W\cong W.
\]
\end{definition}

\begin{lemma}[Shift invariance for weighted approximation]
\label{lem:shift-invariance-rde}
Let
\[
    (A,W,s)
\]
be a weight-cofinal iteration shape.  For every fuzzy domain \(X\) and every
diagram
\[
    D:A\to X,
\]
there is a canonical isomorphism
\[
    W*(Ds)
    \cong
    W*D.
\]
\end{lemma}

\begin{proof}
Apply Proposition~\ref{prop:finite-flat-cofinality} to the \(W\)-cofinal map
\[
    s:(A,W)\to(A,W).
\]
\end{proof}

\section{Finite-flat cofinal bilimits}
\label{sec:finite-flat-cofinal-bilimits}

\begin{definition}[\(W\)-bilimit cone]
\label{def:W-bilimit-cone}
Let
\[
    (A,W)
\]
be a finite-flat approximation shape, and let
\[
    D:A\rightsquigarrow\mathbf{FDom}_{\mathbb V}^{\triangleleft}
\]
be a split \(A\)-indexed projection-pair system.

A \emph{\(W\)-bilimit cone} over \(D\) consists of a fuzzy domain
\[
    D_\infty
\]
and a compatible projection-pair cone
\[
    (\varepsilon_a,\pi_a):D_a\triangleleft D_\infty
    \qquad(a\in A),
\]
such that:

\begin{enumerate}
    \item the cone represents compatible projection-pair cones: for every fuzzy
    domain \(Z\), composition with the cone induces an equivalence
    \[
        \mathbf{FDom}_{\mathbb V}^{\triangleleft}(D_\infty,Z)
        \simeq
        \operatorname{Cone}_{\triangleleft}(D;Z);
    \]

    \item if
    \[
        r:A\to[D_\infty,D_\infty]_{\mathrm{ff}}
    \]
    is the approximant functor
    \[
        r_a=\varepsilon_a\pi_a,
    \]
    then the finite-flat approximation identity holds:
    \[
        W*r
        \cong
        \id_{D_\infty}
    \]
    in
    \[
        [D_\infty,D_\infty]_{\mathrm{ff}}.
    \]
\end{enumerate}
\end{definition}

\begin{lemma}[Uniqueness of \(W\)-bilimits]
\label{lem:uniqueness-W-bilimits}
Any two \(W\)-bilimit cones over the same split projection-pair system are
equivalent, and the equivalence is unique up to invertible compatible \(2\)-cell.
\end{lemma}

\begin{proof}
This is the uniqueness part of the representing equivalence in
Definition~\ref{def:W-bilimit-cone}.  The approximation identity is transported
along the resulting compatible equivalence because it is expressed in the
internal hom of endomorphisms.
\end{proof}

\begin{definition}[Bilimit-cofinal reindexing]
\label{def:bilimit-cofinal-reindexing}
Let
\[
    h:(A,W)\to(B,U)
\]
be a cofinal map of finite-flat approximation shapes, and let
\[
    D:B\rightsquigarrow\mathbf{FDom}_{\mathbb V}^{\triangleleft}
\]
be a split \(B\)-indexed projection-pair system.

We say that \(h\) is \emph{bilimit-cofinal for \(D\)} if restriction along
\(h\) induces, for every fuzzy domain \(Z\), an equivalence of compatible cone
categories
\[
    \operatorname{Cone}_{\triangleleft}(D;Z)
    \simeq
    \operatorname{Cone}_{\triangleleft}(Dh;Z).
\]
Under this equivalence, the derived approximant functor for a reindexed cone is
the precomposition of the original derived approximant functor with \(h\).
\end{definition}

\begin{lemma}[Cofinal reindexing of \(W\)-bilimits]
\label{lem:cofinal-reindexing-W-bilimits}
Let
\[
    h:(A,W)\to(B,U)
\]
be cofinal and bilimit-cofinal for
\[
    D:B\rightsquigarrow\mathbf{FDom}_{\mathbb V}^{\triangleleft}.
\]
Then \(D\) has a \(U\)-bilimit with vertex \(D_\infty\) if and only if the
reindexed system
\[
    Dh:A\rightsquigarrow\mathbf{FDom}_{\mathbb V}^{\triangleleft}
\]
has a \(W\)-bilimit with the corresponding transported cone on the same object
\(D_\infty\).
\end{lemma}

\begin{proof}
The projection-pair universal property is transported along the equivalence of
cone categories.

It remains to compare approximation identities.  Let
\[
    r:B\to[D_\infty,D_\infty]_{\mathrm{ff}}
\]
be the approximant functor for the \(B\)-indexed cone.  The approximant functor
for the reindexed cone is
\[
    rh:A\to[D_\infty,D_\infty]_{\mathrm{ff}}.
\]
Since
\[
    h_!W\cong U,
\]
Proposition~\ref{prop:finite-flat-cofinality} gives
\[
    W*(rh)
    \cong
    U*r.
\]
Thus
\[
    U*r\cong\id_{D_\infty}
\]
if and only if
\[
    W*(rh)\cong\id_{D_\infty}.
\]
\end{proof}

\begin{theorem}[Shifted bilimits]
\label{thm:shifted-bilimits-ff}
Let
\[
    (A,W,s)
\]
be a weight-cofinal iteration shape.  Let
\[
    D:A\rightsquigarrow\mathbf{FDom}_{\mathbb V}^{\triangleleft}
\]
be a split \(A\)-indexed projection-pair system with \(W\)-bilimit
\(D_\infty\).  If
\[
    s:(A,W)\to(A,W)
\]
is bilimit-cofinal for \(D\), then \(D_\infty\) is also a \(W\)-bilimit of the
shifted system
\[
    Ds:A\rightsquigarrow\mathbf{FDom}_{\mathbb V}^{\triangleleft}.
\]
\end{theorem}

\begin{proof}
Apply Lemma~\ref{lem:cofinal-reindexing-W-bilimits} to the map
\[
    s:(A,W)\to(A,W).
\]
Equivalently, if
\[
    r:A\to[D_\infty,D_\infty]_{\mathrm{ff}}
\]
is the approximant functor for the original cone, then the approximant functor
for the shifted cone is
\[
    rs:A\to[D_\infty,D_\infty]_{\mathrm{ff}},
\]
and
\[
    W*(rs)
    \cong
    W*r
    \cong
    \id_{D_\infty}
\]
by Lemma~\ref{lem:shift-invariance-rde}.
\end{proof}

\section{Recursive domain equations by cofinal approximation}
\label{sec:rde-by-cofinal-approximation}

\begin{definition}[\(W\)-bilimit-preserving endopseudofunctor]
\label{def:W-bilimit-preserving-endofunctor}
Let
\[
    (A,W)
\]
be a finite-flat approximation shape.  A projection-pair endopseudofunctor
\[
    F:\mathbf{FDom}_{\mathbb V}^{\triangleleft}
      \to
      \mathbf{FDom}_{\mathbb V}^{\triangleleft}
\]
is \emph{\(W\)-bilimit-preserving} if it sends \(W\)-bilimit cones of split
\(A\)-indexed projection-pair systems to \(W\)-bilimit cones.
\end{definition}

\begin{theorem}[Finite-flat cofinal recursive-domain-equation theorem]
\label{thm:finite-flat-cofinal-rde}
Let
\[
    (A,W,s)
\]
be a weight-cofinal iteration shape.  Let
\[
    F:\mathbf{FDom}_{\mathbb V}^{\triangleleft}
      \to
      \mathbf{FDom}_{\mathbb V}^{\triangleleft}
\]
be a \(W\)-bilimit-preserving projection-pair endopseudofunctor.

Let
\[
    D:A\rightsquigarrow\mathbf{FDom}_{\mathbb V}^{\triangleleft}
\]
be a split \(A\)-indexed projection-pair system equipped with an isomorphism of
split projection-pair systems
\[
    Ds\cong FD.
\]
If \(s\) is bilimit-cofinal for \(D\) and \(D_\infty\) is a \(W\)-bilimit of
\(D\), then there is an equivalence of fuzzy domains
\[
    D_\infty\simeq F(D_\infty).
\]
\end{theorem}

\begin{proof}
By Theorem~\ref{thm:shifted-bilimits-ff}, \(D_\infty\) is a \(W\)-bilimit of
the shifted system \(Ds\).  Since
\[
    Ds\cong FD
\]
as split projection-pair systems, \(D_\infty\) is also a \(W\)-bilimit of
\(FD\).

Since \(F\) preserves \(W\)-bilimits,
\[
    F(D_\infty)
\]
is also a \(W\)-bilimit of \(FD\).  By
Lemma~\ref{lem:uniqueness-W-bilimits},
\[
    D_\infty\simeq F(D_\infty).
\]
\end{proof}

\section{Composition and external-product internal-hom approximation}
\label{sec:composition-and-internal-homs-rde}

\begin{lemma}[Separate continuity of composition]
\label{lem:separate-continuity-composition-rde}
For fuzzy domains \(A,B,C\), composition determines a finite-flat bimorphism
\[
    [B,C]_{\mathrm{ff}}\tensor [A,B]_{\mathrm{ff}}
    \longrightarrow
    [A,C]_{\mathrm{ff}},
    \qquad
    (g,f)\longmapsto gf.
\]
Equivalently, it induces a fuzzy-domain morphism
\[
    [B,C]_{\mathrm{ff}}
    \boxtimes_{\mathrm{ff}}
    [A,B]_{\mathrm{ff}}
    \longrightarrow
    [A,C]_{\mathrm{ff}}.
\]
\end{lemma}

\begin{proof}
Finite-flat colimits in internal homs are computed pointwise.

Fix
\[
    f:A\to B.
\]
For a finite-flat weight
\[
    W:K^{\op}\to\underline{\mathbb V}
\]
and a diagram
\[
    H:K\to[B,C]_{\mathrm{ff}},
\]
one has pointwise
\[
\begin{aligned}
    \bigl(W*H\bigr)\circ f
    &=
    a\longmapsto (W*H)(fa)
    \\
    &\cong
    a\longmapsto W*(k\mapsto Hk(fa))
    \\
    &\cong
    W*(k\mapsto Hk\circ f).
\end{aligned}
\]
Hence precomposition by \(f\) preserves finite-flat colimits.

Fix
\[
    g:B\to C.
\]
Since \(g\) is a fuzzy-domain morphism, it preserves finite-flat colimits in
\(B\).  Therefore, for every finite-flat
\[
    W:K^{\op}\to\underline{\mathbb V}
\]
and every diagram
\[
    H:K\to[A,B]_{\mathrm{ff}},
\]
one has pointwise
\[
    g\bigl(W*(k\mapsto Hk(a))\bigr)
    \cong
    W*(k\mapsto g(Hk(a))).
\]
Thus postcomposition by \(g\) preserves finite-flat colimits.  Therefore
composition is finite-flat-continuous separately in both variables.
\end{proof}

Let
\[
    X:A\rightsquigarrow\mathbf{FDom}_{\mathbb V}^{\triangleleft},
    \qquad
    Y:A\rightsquigarrow\mathbf{FDom}_{\mathbb V}^{\triangleleft}
\]
be split \(A\)-indexed projection-pair systems.  Their associated
\emph{external-product internal-hom system}
\[
    [X,Y]_{\mathrm{ext}}
    :
    A\tensor A
    \rightsquigarrow
    \mathbf{FDom}_{\mathbb V}^{\triangleleft}
\]
is defined on objects by
\[
    [X,Y]_{\mathrm{ext}}(a,a')
    =
    [X_{a'},Y_a]_{\mathrm{ff}}.
\]

For a support arrow
\[
    \xi:(a,a')\to(b,b')
\]
in \((A\tensor A)_0\), let
\[
    \xi_1:a\to b,
    \qquad
    \xi_2:a'\to b'
\]
be its two support projections.  The transition projection pair
\[
    [X_{a'},Y_a]_{\mathrm{ff}}
    \triangleleft
    [X_{b'},Y_b]_{\mathrm{ff}}
\]
has embedding
\[
    \mathsf E_\xi(f)
    =
    e^Y_{\xi_1}\, f\, p^X_{\xi_2},
\]
and projection
\[
    \mathsf P_\xi(g)
    =
    p^Y_{\xi_1}\, g\, e^X_{\xi_2}.
\]

\begin{lemma}[Internal-hom transitions are projection pairs]
\label{lem:internal-hom-transitions-projection-pairs}
For every support arrow \(\xi\) in \((A\tensor A)_0\), the pair
\[
    (\mathsf E_\xi,\mathsf P_\xi)
\]
defined above is a fuzzy projection pair.
\end{lemma}

\begin{proof}
The two maps preserve finite-flat colimits by
Lemma~\ref{lem:separate-continuity-composition-rde}.

For
\[
    f:X_{a'}\to Y_a,
\]
one has
\[
\begin{aligned}
    \mathsf P_\xi\mathsf E_\xi(f)
    &=
    p^Y_{\xi_1}e^Y_{\xi_1}\, f\,
    p^X_{\xi_2}e^X_{\xi_2}
    \\
    &\cong
    f,
\end{aligned}
\]
using the unit isomorphisms of the projection pairs for \(X\) and \(Y\).

For
\[
    g:X_{b'}\to Y_b,
\]
one has
\[
    \mathsf E_\xi\mathsf P_\xi(g)
    =
    e^Y_{\xi_1}p^Y_{\xi_1}\, g\,
    e^X_{\xi_2}p^X_{\xi_2}.
\]
The counit comparisons
\[
    e^Y_{\xi_1}p^Y_{\xi_1}\Rightarrow \id_{Y_b},
    \qquad
    e^X_{\xi_2}p^X_{\xi_2}\Rightarrow \id_{X_{b'}}
\]
induce
\[
    \mathsf E_\xi\mathsf P_\xi(g)
    \Rightarrow
    g.
\]
The triangle identities follow from the triangle identities of the projection
pairs in \(X\) and \(Y\).
\end{proof}

Now suppose
\[
    (\varepsilon_a^X,\pi_a^X):X_a\triangleleft X_\infty,
    \qquad
    (\varepsilon_a^Y,\pi_a^Y):Y_a\triangleleft Y_\infty
\]
are compatible projection-pair cones.  Define
\[
    E_{a,a'}:[X_{a'},Y_a]_{\mathrm{ff}}
    \to [X_\infty,Y_\infty]_{\mathrm{ff}}
\]
and
\[
    \Pi_{a,a'}:[X_\infty,Y_\infty]_{\mathrm{ff}}
    \to [X_{a'},Y_a]_{\mathrm{ff}}
\]
by
\[
    E_{a,a'}(f)=\varepsilon_a^Y f\pi_{a'}^X,
\]
and
\[
    \Pi_{a,a'}(g)=\pi_a^Y g\varepsilon_{a'}^X.
\]

\begin{lemma}[Internal-hom cone components are projection pairs]
\label{lem:internal-hom-cone-components-projection-pairs}
For every \(a,a'\), the pair
\[
    (E_{a,a'},\Pi_{a,a'}):
    [X_{a'},Y_a]_{\mathrm{ff}}
    \triangleleft
    [X_\infty,Y_\infty]_{\mathrm{ff}}
\]
is a fuzzy projection pair.
\end{lemma}

\begin{proof}
The maps \(E_{a,a'}\) and \(\Pi_{a,a'}\) preserve finite-flat colimits by
Lemma~\ref{lem:separate-continuity-composition-rde}.

For \(f:X_{a'}\to Y_a\),
\[
\begin{aligned}
    \Pi_{a,a'}E_{a,a'}(f)
    &=
    \pi_a^Y\varepsilon_a^Y f\pi_{a'}^X\varepsilon_{a'}^X
    \\
    &\cong
    f,
\end{aligned}
\]
using the unit isomorphisms of the projection pairs in the two bilimit cones.

For \(g:X_\infty\to Y_\infty\),
\[
    E_{a,a'}\Pi_{a,a'}(g)
    =
    \varepsilon_a^Y\pi_a^Y g\varepsilon_{a'}^X\pi_{a'}^X.
\]
The counit comparisons
\[
    \varepsilon_a^Y\pi_a^Y\Rightarrow\id_{Y_\infty},
    \qquad
    \varepsilon_{a'}^X\pi_{a'}^X\Rightarrow\id_{X_\infty}
\]
induce
\[
    E_{a,a'}\Pi_{a,a'}(g)\Rightarrow g.
\]
The triangle identities follow from the triangle identities of the two
projection-pair cones.
\end{proof}

\begin{theorem}[External-product internal-hom approximation theorem]
\label{thm:external-product-internal-hom-approximation}
Let
\[
    (A,W)
\]
be a finite-flat approximation shape.  Suppose
\[
    X_\infty
\]
and
\[
    Y_\infty
\]
are \(W\)-bilimits of split \(A\)-indexed projection-pair systems
\[
    X
    \qquad\text{and}\qquad
    Y.
\]
Then the external-product internal-hom cone
\[
    (E_{a,a'},\Pi_{a,a'}):
    [X_{a'},Y_a]_{\mathrm{ff}}
    \triangleleft
    [X_\infty,Y_\infty]_{\mathrm{ff}}
\]
has approximants satisfying
\[
    (W\boxtimes W)*
    \bigl((a,a')\mapsto E_{a,a'}\Pi_{a,a'}\bigr)
    \cong
    \id_{[X_\infty,Y_\infty]_{\mathrm{ff}}}.
\]
\end{theorem}

\begin{proof}
Let
\[
    x_{a'}=\varepsilon_{a'}^X\pi_{a'}^X
    \in [X_\infty,X_\infty]_{\mathrm{ff}},
    \qquad
    y_a=\varepsilon_a^Y\pi_a^Y
    \in [Y_\infty,Y_\infty]_{\mathrm{ff}}.
\]
Then
\[
    E_{a,a'}\Pi_{a,a'}(g)=y_a\,g\,x_{a'}.
\]
By separate finite-flat-continuity of composition and finite-flat Fubini, the
\((W\boxtimes W)\)-weighted colimit of these endofunctors is
\[
    g
    \longmapsto
    \bigl(W*(a\mapsto y_a)\bigr)\,
    g\,
    \bigl(W*(a'\mapsto x_{a'})\bigr).
\]
Since \(X_\infty\) and \(Y_\infty\) are \(W\)-bilimits,
\[
    W*(a'\mapsto x_{a'})
    \cong
    \id_{X_\infty},
    \qquad
    W*(a\mapsto y_a)
    \cong
    \id_{Y_\infty}.
\]
Thus the displayed endofunctor is naturally isomorphic to the identity on
\[
    [X_\infty,Y_\infty]_{\mathrm{ff}}.
\]
\end{proof}

\section{Canonical common-successor comparisons}
\label{sec:canonical-common-successor-comparisons}

The general enriched setting permits canonical comparison maps from double
approximations to single approximations.  These comparisons are not generally
invertible; this is the precise point at which the absence of contraction
matters.

Let
\[
    (P,R)
\]
be common-successor-cofinal.  Let
\[
    H:P\rightsquigarrow\mathbf{FDom}_{\mathbb V}^{\triangleleft}
\]
be a split \(P\)-indexed projection-pair system.  Let
\[
    (E_i,\Pi_i):H_i\triangleleft H_\infty
    \qquad(i\in P)
\]
be a compatible cone, and let
\[
    (\delta_i,\rho_i):H_i\triangleleft Z
    \qquad(i\in P)
\]
be another compatible cone.

Write
\[
    L_i=\delta_i\Pi_i:H_\infty\to Z,
    \qquad
    M_i=E_i\rho_i:Z\to H_\infty.
\]

\begin{lemma}[Canonical tail restrictions]
\label{lem:canonical-tail-restrictions}
Assume \((P,R)\) is tail-cofinal.  For every \(i\in P\),
\[
    R*(j\mapsto L_jE_i)
    \cong
    \delta_i,
\]
and
\[
    R*(j\mapsto \Pi_iM_j)
    \cong
    \rho_i.
\]
\end{lemma}

\begin{proof}
Let
\[
    \tau_i:\operatorname{Tail}_P(i)\to P
\]
be the canonical support tail.  Since \((P,R)\) is tail-cofinal,
\[
    (\tau_i)_!R^i\cong R.
\]
Therefore, by Proposition~\ref{prop:finite-flat-cofinality},
\[
    R*(j\mapsto L_jE_i)
    \cong
    R^i*((j,u)\mapsto L_jE_i).
\]

For an object \((j,u:i\to j)\) of \(\operatorname{Tail}_P(i)\), cone
compatibility gives
\[
    E_j e_u\cong E_i,
    \qquad
    p_u\Pi_j\cong \Pi_i,
    \qquad
    \delta_j e_u\cong \delta_i.
\]
Hence
\[
\begin{aligned}
    L_jE_i
    &=
    \delta_j\Pi_jE_i
    \\
    &\cong
    \delta_j\Pi_jE_j e_u
    \\
    &\cong
    \delta_j e_u
    \\
    &\cong
    \delta_i.
\end{aligned}
\]
Thus the restricted diagram is constant at \(\delta_i\), and
\[
    R^i*((j,u)\mapsto L_jE_i)
    \cong
    R^i*\Delta(\delta_i)
    \cong
    \delta_i
\]
by Lemma~\ref{lem:constant-diagrams-canonical-tail}.

The proof of
\[
    R*(j\mapsto \Pi_iM_j)
    \cong
    \rho_i
\]
is analogous.  For \((j,u:i\to j)\), cone compatibility gives
\[
    p_u\rho_j\cong \rho_i
\]
and
\[
\begin{aligned}
    \Pi_iM_j
    &=
    \Pi_iE_j\rho_j
    \\
    &\cong
    p_u\Pi_jE_j\rho_j
    \\
    &\cong
    p_u\rho_j
    \\
    &\cong
    \rho_i.
\end{aligned}
\]
\end{proof}

\begin{lemma}[Canonical double-to-single comparison]
\label{lem:canonical-double-to-single-comparison}
Assume \((P,R)\) is common-successor-cofinal.  Then there is a canonical
comparison
\[
    (R\boxtimes R)*
    \bigl((i,j)\mapsto M_iL_j\bigr)
    \Rightarrow
    R*(k\mapsto E_k\Pi_k).
\]
Moreover, there is a canonical counit comparison
\[
    (R\boxtimes R)*
    \bigl((i,j)\mapsto L_iM_j\bigr)
    \Rightarrow
    \id_Z.
\]
\end{lemma}

\begin{proof}
Let
\[
    R^{\mathrm{com}}=m_P^*R
\]
be the canonical common-successor weight.  By common-successor-cofinality,
\[
    (c_P)_!R^{\mathrm{com}}\cong R\boxtimes R,
    \qquad
    (m_P)_!R^{\mathrm{com}}\cong R.
\]
Hence
\[
\begin{aligned}
    (R\boxtimes R)*
    \bigl((i,j)\mapsto M_iL_j\bigr)
    &\cong
    R^{\mathrm{com}}*
    \bigl((i,j,k,u,v)\mapsto M_iL_j\bigr).
\end{aligned}
\]

For an object
\[
    (i,j,k,u:i\to k,v:j\to k)
\]
of \(\operatorname{Com}(P)\), cone compatibility gives
\[
    E_k e_u\cong E_i,
    \qquad
    p_u\rho_k\cong \rho_i,
\]
and
\[
    \delta_k e_v\cong \delta_j,
    \qquad
    p_v\Pi_k\cong \Pi_j.
\]
Therefore
\[
\begin{aligned}
    M_iL_j
    &=
    E_i\rho_i\delta_j\Pi_j
    \\
    &\cong
    E_k e_u p_u\rho_k\delta_k e_vp_v\Pi_k
    \\
    &\cong
    E_k e_u p_u e_vp_v\Pi_k
    \\
    &\Rightarrow
    E_k\Pi_k.
\end{aligned}
\]
The final comparison is induced by the counits
\[
    e_up_u\Rightarrow \id_{H_k},
    \qquad
    e_vp_v\Rightarrow \id_{H_k}.
\]
These comparisons are coherent in \(\operatorname{Com}(P)\).  Taking the
\(R^{\mathrm{com}}\)-weighted colimit and using
\[
    (m_P)_!R^{\mathrm{com}}\cong R
\]
gives the first comparison.

For the second comparison, again use cofinality of \(c_P\).  For an object
\[
    (i,j,k,u:i\to k,v:j\to k)
\]
there is a canonical comparison
\[
    L_iM_j\Rightarrow \id_Z.
\]
Indeed,
\[
\begin{aligned}
    L_iM_j
    &=
    \delta_i\Pi_iE_j\rho_j
    \\
    &\cong
    \delta_k e_u p_u\Pi_kE_k e_vp_v\rho_k
    \\
    &\Rightarrow
    \delta_k\rho_k
    \\
    &\Rightarrow
    \id_Z.
\end{aligned}
\]
The final comparison is the counit of
\[
    (\delta_k,\rho_k):H_k\triangleleft Z.
\]
Taking the \(R^{\mathrm{com}}\)-weighted colimit and using
Lemma~\ref{lem:constant-diagrams-finite-flat} for the constant diagram at
\(\id_Z\) gives the second comparison.
\end{proof}

\begin{theorem}[Lax reconstruction of external-product cones]
\label{thm:lax-reconstruction-internal-hom-cones}
Let
\[
    (A,W)
\]
be linear-cofinal.  Put
\[
    P=A\tensor A,
    \qquad
    R=W\boxtimes W.
\]
Let
\[
    H_\infty=[X_\infty,Y_\infty]_{\mathrm{ff}}
\]
with external-product approximants
\[
    (E_i,\Pi_i):H_i\triangleleft H_\infty
    \qquad(i\in P).
\]
Let
\[
    (\delta_i,\rho_i):H_i\triangleleft Z
\]
be a compatible cone.  Define
\[
    \delta
    =
    R*(i\mapsto \delta_i\Pi_i):
    H_\infty\to Z
\]
and
\[
    \rho
    =
    R*(i\mapsto E_i\rho_i):
    Z\to H_\infty.
\]
Then
\[
    \delta E_i\cong \delta_i,
    \qquad
    \Pi_i\rho\cong \rho_i,
\]
and there are canonical comparisons
\[
    \rho\delta\Rightarrow \id_{H_\infty},
    \qquad
    \delta\rho\Rightarrow \id_Z.
\]
\end{theorem}

\begin{proof}
The restriction equations follow from
Lemma~\ref{lem:canonical-tail-restrictions}.

For the first comparison, separate finite-flat-continuity of composition and
finite-flat Fubini give
\[
    \rho\delta
    \cong
    (R\boxtimes R)*
    \bigl((i,j)\mapsto M_iL_j\bigr).
\]
By Lemma~\ref{lem:canonical-double-to-single-comparison},
\[
    (R\boxtimes R)*
    \bigl((i,j)\mapsto M_iL_j\bigr)
    \Rightarrow
    R*(k\mapsto E_k\Pi_k).
\]
By Theorem~\ref{thm:external-product-internal-hom-approximation},
\[
    R*(k\mapsto E_k\Pi_k)
    \cong
    \id_{H_\infty}.
\]
Thus
\[
    \rho\delta\Rightarrow\id_{H_\infty}.
\]
The comparison
\[
    \delta\rho\Rightarrow \id_Z
\]
is the second comparison of
Lemma~\ref{lem:canonical-double-to-single-comparison}.
\end{proof}

\section{Linear application in the general finite-flat setting}
\label{sec:linear-application-general-ff}

The preceding theorem reconstructs canonical lax application data from
external-product finite-stage cones.  It does not produce a projection-pair
retract in full generality; the missing ingredient is precisely a diagonal or
contraction comparison.

\begin{definition}[Linear reflexive fuzzy domain]
\label{def:linear-reflexive-fuzzy-domain-ff}
A \emph{linear reflexive fuzzy domain} is a fuzzy domain \(D\) equipped with a
fuzzy projection pair
\[
    [D,D]_{\mathrm{ff}}\triangleleft D.
\]
Equivalently, it consists of fuzzy-domain morphisms
\[
    \operatorname{encode}:[D,D]_{\mathrm{ff}}\to D,
    \qquad
    \operatorname{decode}:D\to [D,D]_{\mathrm{ff}},
\]
together with an adjunction
\[
    \operatorname{encode}\dashv \operatorname{decode}
\]
whose unit is invertible.  Thus
\[
    \operatorname{decode}\operatorname{encode}
    \cong
    \id_{[D,D]_{\mathrm{ff}}},
\]
and there is a counit comparison
\[
    \operatorname{encode}\operatorname{decode}
    \Rightarrow
    \id_D
\]
satisfying the triangle identities.

It is \emph{extensional} if the counit is also invertible.
\end{definition}

The associated application morphism is
\[
    D\boxtimes_{\mathrm{ff}}D
    \xrightarrow{\operatorname{decode}\boxtimes 1_D}
    [D,D]_{\mathrm{ff}}\boxtimes_{\mathrm{ff}}D
    \xrightarrow{\operatorname{ev}}
    D.
\]

\begin{definition}[\((G,T)\)-linear reflexive fuzzy domain]
\label{def:mixed-linear-reflexive-fuzzy-domain-ff}
Let
\[
    G,T:\mathbf{FDom}_{\mathbb V}^{\triangleleft}
      \to
      \mathbf{FDom}_{\mathbb V}^{\triangleleft}
\]
be projection-pair endopseudofunctors.  A
\emph{\((G,T)\)-linear reflexive fuzzy domain} is a fuzzy domain \(D\) equipped
with a fuzzy projection pair
\[
    [GD,TD]_{\mathrm{ff}}\triangleleft D.
\]
The associated application morphism has type
\[
    D\boxtimes_{\mathrm{ff}}GD\to TD
\]
and is the composite
\[
    D\boxtimes_{\mathrm{ff}}GD
    \xrightarrow{\operatorname{decode}\boxtimes 1_{GD}}
    [GD,TD]_{\mathrm{ff}}\boxtimes_{\mathrm{ff}}GD
    \xrightarrow{\operatorname{ev}}
    TD.
\]
\end{definition}

\begin{definition}[External finite-stage endomorphism system]
\label{def:external-finite-stage-endomorphism-system}
Let
\[
    D:A\rightsquigarrow\mathbf{FDom}_{\mathbb V}^{\triangleleft}
\]
be a split \(A\)-indexed projection-pair system.  Its
\emph{external finite-stage endomorphism system} is the split
\((A\tensor A)\)-indexed projection-pair system
\[
    \operatorname{End}_{\mathrm{ext}}(D)
    :
    A\tensor A
    \rightsquigarrow
    \mathbf{FDom}_{\mathbb V}^{\triangleleft}
\]
defined by
\[
    \operatorname{End}_{\mathrm{ext}}(D)(a,a')
    =
    [D_{a'},D_a]_{\mathrm{ff}}.
\]
\end{definition}

\begin{definition}[Linear self-application cone]
\label{def:linear-self-application-cone}
Let
\[
    (A,W)
\]
be linear-cofinal, and let
\[
    D:A\rightsquigarrow\mathbf{FDom}_{\mathbb V}^{\triangleleft}
\]
be a split \(A\)-indexed projection-pair system with \(W\)-bilimit
\[
    D_\infty.
\]
A \emph{linear self-application cone} on \(D\) is a compatible projection-pair
cone
\[
    (\kappa_{a,a'},\lambda_{a,a'}):
    [D_{a'},D_a]_{\mathrm{ff}}
    \triangleleft
    D_\infty
\]
over the external finite-stage endomorphism system
\[
    \operatorname{End}_{\mathrm{ext}}(D).
\]
\end{definition}

\begin{theorem}[Pure lax linear application theorem]
\label{thm:pure-lax-linear-application}
Let
\[
    (A,W)
\]
be linear-cofinal, and suppose
\[
    D_\infty
\]
is a \(W\)-bilimit of a split \(A\)-indexed projection-pair system \(D\).
Let
\[
    H_\infty=[D_\infty,D_\infty]_{\mathrm{ff}}.
\]
If \(D\) carries a linear self-application cone, then there are canonical
fuzzy-domain morphisms
\[
    \operatorname{encode}:H_\infty\to D_\infty,
    \qquad
    \operatorname{decode}:D_\infty\to H_\infty
\]
defined by
\[
    \operatorname{encode}
    =
    (W\boxtimes W)*
    \bigl((a,a')\mapsto
        \kappa_{a,a'}\Pi_{a,a'}
    \bigr),
\]
and
\[
    \operatorname{decode}
    =
    (W\boxtimes W)*
    \bigl((a,a')\mapsto
        E_{a,a'}\lambda_{a,a'}
    \bigr).
\]
They satisfy canonical comparisons
\[
    \operatorname{decode}\operatorname{encode}
    \Rightarrow
    \id_{H_\infty},
    \qquad
    \operatorname{encode}\operatorname{decode}
    \Rightarrow
    \id_{D_\infty}.
\]
Consequently there is a canonical resource-sensitive application morphism
\[
    D_\infty\boxtimes_{\mathrm{ff}}D_\infty
    \longrightarrow
    D_\infty.
\]
\end{theorem}

\begin{proof}
Apply Theorem~\ref{thm:lax-reconstruction-internal-hom-cones} to the external
endomorphism system
\[
    (a,a')\longmapsto [D_{a'},D_a]_{\mathrm{ff}}
\]
and to the compatible cone
\[
    (\kappa_{a,a'},\lambda_{a,a'}):
    [D_{a'},D_a]_{\mathrm{ff}}
    \triangleleft
    D_\infty.
\]
The resulting maps are precisely the displayed
\(\operatorname{encode}\) and \(\operatorname{decode}\).  The application
morphism is
\[
    D_\infty\boxtimes_{\mathrm{ff}}D_\infty
    \xrightarrow{\operatorname{decode}\boxtimes 1}
    [D_\infty,D_\infty]_{\mathrm{ff}}
    \boxtimes_{\mathrm{ff}}
    D_\infty
    \xrightarrow{\operatorname{ev}}
    D_\infty.
\]
\end{proof}

\begin{definition}[Mixed external finite-stage function-space system]
\label{def:mixed-external-finite-stage-function-space-system}
Let
\[
    G,T:\mathbf{FDom}_{\mathbb V}^{\triangleleft}
      \to
      \mathbf{FDom}_{\mathbb V}^{\triangleleft}
\]
be projection-pair endopseudofunctors, and let
\[
    D:A\rightsquigarrow\mathbf{FDom}_{\mathbb V}^{\triangleleft}
\]
be a split \(A\)-indexed projection-pair system.  The
\emph{mixed external finite-stage function-space system} is the split
\((A\tensor A)\)-indexed system
\[
    (a,a')\longmapsto [GD_{a'},TD_a]_{\mathrm{ff}}.
\]
\end{definition}

\begin{definition}[Mixed linear self-application cone]
\label{def:mixed-linear-self-application-cone}
Let
\[
    (A,W)
\]
be linear-cofinal.  Let
\[
    G,T:\mathbf{FDom}_{\mathbb V}^{\triangleleft}
      \to
      \mathbf{FDom}_{\mathbb V}^{\triangleleft}
\]
be projection-pair endopseudofunctors preserving \(W\)-bilimits.  Let
\[
    D:A\rightsquigarrow\mathbf{FDom}_{\mathbb V}^{\triangleleft}
\]
be a split \(A\)-indexed projection-pair system with \(W\)-bilimit
\[
    D_\infty.
\]
A \emph{mixed linear self-application cone} is a compatible projection-pair cone
\[
    [GD_{a'},TD_a]_{\mathrm{ff}}
    \triangleleft
    D_\infty
\]
over the mixed external finite-stage function-space system.
\end{definition}

\begin{theorem}[Mixed lax linear application theorem]
\label{thm:mixed-lax-linear-application}
Let
\[
    (A,W)
\]
be linear-cofinal.  Let
\[
    G,T:\mathbf{FDom}_{\mathbb V}^{\triangleleft}
      \to
      \mathbf{FDom}_{\mathbb V}^{\triangleleft}
\]
be projection-pair endopseudofunctors preserving \(W\)-bilimits.  Suppose
\[
    D_\infty
\]
is a \(W\)-bilimit of a split \(A\)-indexed projection-pair system \(D\).
If \(D\) carries a mixed linear self-application cone, then there are canonical
morphisms
\[
    \operatorname{encode}:[GD_\infty,TD_\infty]_{\mathrm{ff}}\to D_\infty,
\]
and
\[
    \operatorname{decode}:D_\infty\to[GD_\infty,TD_\infty]_{\mathrm{ff}},
\]
with canonical comparisons
\[
    \operatorname{decode}\operatorname{encode}
    \Rightarrow
    \id_{[GD_\infty,TD_\infty]_{\mathrm{ff}}},
    \qquad
    \operatorname{encode}\operatorname{decode}
    \Rightarrow
    \id_{D_\infty}.
\]
The associated application morphism has type
\[
    D_\infty\boxtimes_{\mathrm{ff}}GD_\infty\to TD_\infty.
\]
\end{theorem}

\begin{proof}
Since \(G\) and \(T\) preserve \(W\)-bilimits,
\[
    GD_\infty
\]
is a \(W\)-bilimit of \(GD\), and
\[
    TD_\infty
\]
is a \(W\)-bilimit of \(TD\).  Apply
Theorem~\ref{thm:lax-reconstruction-internal-hom-cones} to the mixed external
finite-stage function-space system.
\end{proof}

\subsection{Finite-stage generation of self-application cones}

The preceding lax application theorem requires a compatible cone from the
finite-stage function-space system into the bilimit.  The following construction
supplies such data over the canonical common-successor category.

\begin{definition}[Finite-stage linear self-application data]
\label{def:finite-stage-linear-self-application-data}
Let
\[
    (A,W)
\]
be linear-cofinal, and let
\[
    D:A\rightsquigarrow\mathbf{FDom}_{\mathbb V}^{\triangleleft}
\]
be a split \(A\)-indexed projection-pair system.  Finite-stage linear
self-application data consist of:

\begin{enumerate}
    \item a support endofunctor
    \[
        s:A\to A;
    \]

    \item for every \(k\in A\), a fuzzy projection pair
    \[
        (\sigma_k,\tau_k):
        [D_k,D_k]_{\mathrm{ff}}
        \triangleleft
        D_{sk};
    \]

    \item coherence isomorphisms saying that, for every support arrow
    \[
        w:k\to k',
    \]
    the square of projection pairs comparing
    \[
        [D_k,D_k]_{\mathrm{ff}}
        \triangleleft
        D_{sk}
    \]
    with
    \[
        [D_{k'},D_{k'}]_{\mathrm{ff}}
        \triangleleft
        D_{sk'}
    \]
    is compatible with the internal-hom transition maps and with the transition
    projection pair
    \[
        D_{sk}\triangleleft D_{sk'}.
    \]
\end{enumerate}
\end{definition}

\begin{construction}[Common-successor finite-stage encoding]
\label{con:common-successor-finite-stage-encoding}
Let
\[
    (i,j,k,u:i\to k,v:j\to k)
\]
be an object of \(\operatorname{Com}(A)\).  The transition projection pairs in
\(D\) induce a projection pair
\[
    [D_j,D_i]_{\mathrm{ff}}
    \triangleleft
    [D_k,D_k]_{\mathrm{ff}}.
\]
Its embedding sends
\[
    f:D_j\to D_i
\]
to
\[
    e_u f p_v:D_k\to D_k,
\]
and its projection sends
\[
    g:D_k\to D_k
\]
to
\[
    p_u g e_v:D_j\to D_i.
\]
Composing with
\[
    [D_k,D_k]_{\mathrm{ff}}
    \triangleleft
    D_{sk}
\]
and then with the bilimit cone component
\[
    D_{sk}\triangleleft D_\infty
\]
gives a projection pair
\[
    [D_j,D_i]_{\mathrm{ff}}
    \triangleleft
    D_\infty.
\]
The coherence axioms for finite-stage self-application data make these
projection pairs compatible over \(\operatorname{Com}(A)\).
\end{construction}

\begin{theorem}[Common-successor finite-stage encodings over \(\operatorname{Com}(A)\)]
\label{thm:common-successor-encodings-over-com}
Let
\[
    (A,W)
\]
be linear-cofinal, and let
\[
    D:A\rightsquigarrow\mathbf{FDom}_{\mathbb V}^{\triangleleft}
\]
be a split \(A\)-indexed system with \(W\)-bilimit \(D_\infty\).  Suppose \(D\)
is equipped with finite-stage linear self-application data.  Then the
construction above gives a compatible projection-pair cone over the pullback of
the external endomorphism system along
\[
    c_A:\operatorname{Com}(A)\to A\tensor A.
\]
\end{theorem}

\begin{proof}
For each object
\[
    (i,j,k,u,v)
\]
of \(\operatorname{Com}(A)\), Construction~\ref{con:common-successor-finite-stage-encoding}
gives a projection pair
\[
    [D_j,D_i]_{\mathrm{ff}}\triangleleft D_\infty.
\]
Compatibility with arrows of \(\operatorname{Com}(A)\) follows from the
coherence of the split system \(D\), the coherence of the finite-stage
self-application data, and the calculus of mates.
\end{proof}

\section{Constant-map seeds}
\label{sec:constant-map-seeds-ff}

The following construction supplies standard finite-stage seeds.

\begin{lemma}[Constant finite-flat-continuous maps]
\label{lem:constant-ff-continuous}
Let \(X\) and \(Y\) be fuzzy domains, and let \(y\in Y\).  Since
\(\mathbb V\) is affine, \(y\) determines a constant \(\mathbb V\)-functor
\[
    \operatorname{const}_y:X\to Y.
\]
This constant functor preserves finite-flat weighted colimits.
\end{lemma}

\begin{proof}
The constant assignment is a \(\mathbb V\)-functor because, for every
\(x,x'\in X\), affinity gives a canonical morphism
\[
    X(x,x')\to I,
\]
and the identity of \(y\) gives
\[
    I\to Y(y,y).
\]
Their composite supplies the hom-map
\[
    X(x,x')\to Y(y,y).
\]

Let
\[
    W:K^{\op}\to\underline{\mathbb V}
\]
be finite-flat and let
\[
    D:K\to X
\]
be a diagram.  Since \(W\) is finite-flat,
Lemma~\ref{lem:constant-diagrams-finite-flat} gives
\[
    W*(\operatorname{const}_yD)
    \cong
    y.
\]
Therefore
\[
    \operatorname{const}_y(W*D)
    =
    y
    \cong
    W*(\operatorname{const}_yD).
\]
Thus \(\operatorname{const}_y\) preserves finite-flat weighted colimits.
\end{proof}

\begin{proposition}[Pure constant-map seed]
\label{prop:pure-constant-map-seed-ff}
Suppose a fuzzy domain \(D_0\) has an enriched initial object
\[
    \bot\in D_0.
\]
Define
\[
    e:D_0\to [D_0,D_0]_{\mathrm{ff}}
\]
by
\[
    e(a)=\operatorname{const}_a,
\]
and define
\[
    p:[D_0,D_0]_{\mathrm{ff}}\to D_0
\]
by
\[
    p(f)=f(\bot).
\]
Then
\[
    (e,p):D_0\triangleleft [D_0,D_0]_{\mathrm{ff}}
\]
is a fuzzy projection pair.
\end{proposition}

\begin{proof}
By Lemma~\ref{lem:constant-ff-continuous}, each \(\operatorname{const}_a\)
lies in \([D_0,D_0]_{\mathrm{ff}}\).  The assignment
\[
    a\longmapsto\operatorname{const}_a
\]
is a \(\mathbb V\)-functor because the end defining the internal hom admits the
evident pointwise wedge
\[
    D_0(a,a')\to D_0(a,a').
\]
It preserves finite-flat colimits by Lemma~\ref{lem:constant-ff-continuous} and
Lemma~\ref{lem:constant-diagrams-finite-flat}.

The functor \(p\) is evaluation at \(\bot\).  Evaluation is a
\(\mathbb V\)-functor and preserves pointwise finite-flat colimits.  Thus \(p\)
is a fuzzy-domain morphism.

Moreover,
\[
    p(e(a))
    =
    \operatorname{const}_a(\bot)
    =
    a.
\]
Thus
\[
    pe=\id_{D_0}.
\]

For \(f:D_0\to D_0\), the composite \(ep(f)\) is
\[
    \operatorname{const}_{f(\bot)}.
\]
The enriched initiality of \(\bot\) gives a canonical enriched natural
transformation
\[
    \Delta\bot\Rightarrow 1_{D_0}.
\]
Applying \(f\) yields
\[
    \Delta f(\bot)\Rightarrow f.
\]
This is precisely
\[
    ep(f)\Rightarrow f.
\]
The triangle identities follow from \(pe=\id_{D_0}\) and from functoriality of
the canonical initial-object comparison.  Therefore \(e\dashv p\) with
invertible unit.
\end{proof}

\begin{proposition}[Mixed constant-map seed]
\label{prop:mixed-constant-map-seed-ff}
Let
\[
    G,T:\mathbf{FDom}_{\mathbb V}^{\triangleleft}
      \to
      \mathbf{FDom}_{\mathbb V}^{\triangleleft}
\]
be projection-pair endopseudofunctors.  Suppose \(GD_0\) has an enriched initial
object
\[
    \bot_{GD_0}\in GD_0,
\]
and suppose there is a fuzzy projection pair
\[
    (u,v):D_0\triangleleft TD_0.
\]
Define
\[
    e:D_0\to[GD_0,TD_0]_{\mathrm{ff}}
\]
by
\[
    e(a)=\operatorname{const}_{u(a)},
\]
and define
\[
    p:[GD_0,TD_0]_{\mathrm{ff}}\to D_0
\]
by
\[
    p(f)=v(f(\bot_{GD_0})).
\]
Then
\[
    (e,p):D_0\triangleleft[GD_0,TD_0]_{\mathrm{ff}}
\]
is a fuzzy projection pair.
\end{proposition}

\begin{proof}
Constant maps preserve finite-flat colimits by
Lemma~\ref{lem:constant-ff-continuous}.  Since \(u\) preserves finite-flat
colimits, the functor
\[
    e(a)=\operatorname{const}_{u(a)}
\]
preserves finite-flat colimits.

The functor \(p\) is the composite of evaluation at \(\bot_{GD_0}\) with
\[
    v:TD_0\to D_0.
\]
Evaluation preserves pointwise finite-flat colimits, and \(v\) is a
fuzzy-domain morphism.  Hence \(p\) is a fuzzy-domain morphism.

For \(a\in D_0\),
\[
    p(e(a))
    =
    v(u(a))
    \cong
    a
\]
by the unit isomorphism of the projection pair \((u,v)\).

For \(f:GD_0\to TD_0\), the composite \(ep(f)\) is the constant map at
\[
    u(v(f(\bot_{GD_0}))).
\]
The counit comparison
\[
    uv\Rightarrow \id_{TD_0}
\]
gives
\[
    u(v(f(\bot_{GD_0})))\to f(\bot_{GD_0}).
\]
The enriched initiality of \(\bot_{GD_0}\) gives
\[
    \Delta\bot_{GD_0}\Rightarrow 1_{GD_0}.
\]
Applying \(f\) gives
\[
    \Delta f(\bot_{GD_0})\Rightarrow f.
\]
Composing these transformations gives
\[
    ep(f)\Rightarrow f.
\]
The triangle identities follow from the triangle identities of \((u,v)\) and
from functoriality of the canonical initial-object comparison.
\end{proof}

\section{The posetal reflection}
\label{sec:posetal-reflection-ff-rde}

Let
\[
    \Lambda=(L,\leq,\meet,\join,\tensor,\multimap,\bot_L,\top_L)
\]
be a fuzzy truth-value object.  In the posetal reflection, finite-flat
approximation shapes are fuzzy ideal weights
\[
    W:A^{\op}\to\Lambda.
\]
For a \(\Lambda\)-functor
\[
    h:A\to B,
\]
the lower direct image is computed by
\[
    (h_!W)(b)
    =
    \bigjoin_{a\in A}
        W(a)\tensor B(b,ha).
\]
Cofinality means
\[
    h_!W\cong U.
\]
Since the truth-value fuzzy preorder \(\Lambda\) is separated, the presheaf
fuzzy preorder
\[
    [B^{\op},\Lambda]
\]
is separated.  Hence an isomorphism of weights may be written as equality:
\[
    h_!W=U.
\]

For every finite-flat weight \(W\), the terminal direct image satisfies
\[
    !_!W\cong \top_L.
\]
Equivalently,
\[
    \bigjoin_{a\in A}W(a)=\top_L.
\]
This is the posetal reflection of preservation of the terminal finite-limit
predicate by the character \(\chi_W\).

A projection pair
\[
    (e,p):X\triangleleft Y
\]
is a pair of fuzzy-domain morphisms satisfying
\[
    pe\cong \id_X,
    \qquad
    ep\leq \id_Y
\]
in the induced crisp preorder on fuzzy-domain morphisms, where
\[
    f\leq g
    \quad:\Longleftrightarrow\quad
    [X,Y]_{\mathrm{ff}}(f,g)=\top_L.
\]

A \(W\)-bilimit cone has approximation identity
\[
    W*(a\mapsto \varepsilon_a\pi_a)
    \cong
    \id_{D_\infty}
\]
inside
\[
    [D_\infty,D_\infty]_{\mathrm{ff}}.
\]
The finite-flat cofinal recursive-domain-equation theorem therefore gives:
if
\[
    Ds\cong FD,
\]
if \(s\) is bilimit-cofinal for \(D\), and if \(F\) preserves the corresponding
\(W\)-bilimit, then
\[
    D_\infty\simeq F(D_\infty).
\]

The general enriched internal-hom theorem gives external-product approximation
\[
    (W\boxtimes W)*
    \bigl((a,a')\mapsto E_{a,a'}\Pi_{a,a'}\bigr)
    \cong
    \id_{[X_\infty,Y_\infty]_{\mathrm{ff}}}.
\]
It does not generally collapse to a same-index equation, because there is no
canonical diagonal
\[
    A\to A\tensor A.
\]

\section{Pointed finite-flat minimal invariants}
\label{sec:pointed-ff-minimal-invariants}

The preceding recursive-domain-equation theorem produces solutions from
finite-flat bilimits.  We now isolate the pointed form of the construction,
which is the finite-flat analogue of the classical pointed-dcpo minimal
solution theorem.

The point of pointedness is not that every solution of
\[
    X\simeq F(X)
\]
is unique.  That statement is false already in ordinary pointed dcpos: for
\(F=\id\), every pointed dcpo is a solution.  Rather, pointedness supplies a
canonical initial approximation process.  The resulting bilimit is the
canonical generated invariant of \(F\).

\begin{definition}[Pointed fuzzy domain]
\label{def:pointed-fuzzy-domain}
A \emph{pointed fuzzy domain} is a fuzzy domain \(X\) equipped with an enriched
initial object
\[
    \bot_X\in X.
\]
Thus
\[
    X(\bot_X,x)\cong I
\]
naturally in \(x\in X\).

A morphism of pointed fuzzy domains is a fuzzy-domain morphism
\[
    f:X\to Y.
\]
It is called \emph{strict} if
\[
    f(\bot_X)\cong \bot_Y.
\]
\end{definition}

\begin{lemma}[Projection pairs between pointed domains are strict]
\label{lem:projection-pairs-between-pointed-domains-are-strict}
Let
\[
    (e,p):X\triangleleft Y
\]
be a fuzzy projection pair between pointed fuzzy domains.  Then both
\[
    e:X\to Y
    \qquad
    p:Y\to X
\]
preserve the chosen initial objects:
\[
    e(\bot_X)\cong \bot_Y,
    \qquad
    p(\bot_Y)\cong \bot_X.
\]
\end{lemma}

\begin{proof}
Since \(e\dashv p\), for every \(y\in Y\) there are natural isomorphisms
\[
    Y(e\bot_X,y)
    \cong
    X(\bot_X,py)
    \cong
    I.
\]
Thus \(e\bot_X\) is an enriched initial object of \(Y\), so
\[
    e(\bot_X)\cong\bot_Y.
\]
Using the invertible unit of the projection pair,
\[
    pe\cong\id_X,
\]
we obtain
\[
    p(\bot_Y)
    \cong
    p(e\bot_X)
    \cong
    \bot_X.
\]
\end{proof}

\begin{definition}[Pointed projection-pair category]
\label{def:pointed-projection-pair-category}
Let
\[
    \mathbf{FDom}^{\triangleleft}_{\mathbb V,\bot}
\]
denote the locally ordinary \(2\)-category whose objects are pointed fuzzy
domains, whose \(1\)-cells are fuzzy projection pairs between pointed fuzzy
domains, and whose \(2\)-cells are mate-compatible pairs of enriched natural
transformations.

By Lemma~\ref{lem:projection-pairs-between-pointed-domains-are-strict}, every
\(1\)-cell in
\[
    \mathbf{FDom}^{\triangleleft}_{\mathbb V,\bot}
\]
is automatically strict on the chosen initial objects.
\end{definition}

\begin{definition}[The one-point pointed fuzzy domain]
\label{def:one-point-pointed-fuzzy-domain}
Let
\[
    \mathbf 1_\bot
\]
denote the one-object fuzzy domain.  Its unique object is written
\[
    \bot.
\]
It is pointed by its unique object.
\end{definition}

\begin{lemma}[The one-point domain embeds into every pointed fuzzy domain]
\label{lem:one-point-domain-embeds-into-pointed-domain}
For every pointed fuzzy domain \(X\), there is a canonical fuzzy projection pair
\[
    \mathbf 1_\bot\triangleleft X.
\]
Its embedding sends the unique object of \(\mathbf 1_\bot\) to
\[
    \bot_X,
\]
and its projection is the unique functor
\[
    X\to\mathbf 1_\bot.
\]
\end{lemma}

\begin{proof}
Let
\[
    e:\mathbf 1_\bot\to X
\]
send the unique object to \(\bot_X\), and let
\[
    p:X\to\mathbf 1_\bot
\]
be the unique functor.  Since \(\bot_X\) is enriched initial, for every
\(x\in X\) one has
\[
    X(e*,x)=X(\bot_X,x)\cong I
    \cong
    \mathbf 1_\bot(*,px).
\]
Thus
\[
    e\dashv p.
\]
The unit is the identity, and the counit is the enriched natural transformation
\[
    \Delta\bot_X\Rightarrow\id_X
\]
induced by initiality of \(\bot_X\).
\end{proof}

\begin{definition}[Sequential finite-flat approximation shape]
\label{def:sequential-ff-approximation-shape}
A \emph{sequential finite-flat approximation shape} is a linear-cofinal
finite-flat approximation shape
\[
    (\omega_{\mathbb V},W_\omega)
\]
whose support category is the ordinary preorder \(\omega\), together with a
\(\mathbb V\)-endofunctor
\[
    s:\omega_{\mathbb V}\to\omega_{\mathbb V}
\]
whose action on support objects is
\[
    n\longmapsto n+1,
\]
such that
\[
    s_!W_\omega\cong W_\omega.
\]
\end{definition}

\begin{definition}[Sequential finite-flat cofinality]
\label{def:sequential-ff-cofinality}
The fuzzy cosmos \(\mathbb V\) is \emph{sequentially finite-flat cofinal} if
there is a sequential finite-flat approximation shape
\[
    (\omega_{\mathbb V},W_\omega,s)
\]
such that:

\begin{enumerate}
    \item \(s\) is \(W_\omega\)-cofinal:
    \[
        s_!W_\omega\cong W_\omega;
    \]

    \item \(s\) is bilimit-cofinal for the pointed sequential projection-pair
    systems considered below;

    \item the canonical tails and canonical external-product common-successor
    comparisons are cofinal in the senses of
    Definitions~\ref{def:tail-cofinal-approximation-shape},
    \ref{def:common-successor-cofinal-shape}, and
    \ref{def:linear-cofinal-ff-approximation-shape}.
\end{enumerate}
\end{definition}

\begin{definition}[Pointed projection-pair endopseudofunctor]
\label{def:pointed-projection-pair-endopseudofunctor}
A \emph{pointed projection-pair endopseudofunctor} is a pseudofunctor
\[
    F:
    \mathbf{FDom}^{\triangleleft}_{\mathbb V,\bot}
    \longrightarrow
    \mathbf{FDom}^{\triangleleft}_{\mathbb V,\bot}.
\]
\end{definition}

\begin{definition}[Pointed initial iteration]
\label{def:pointed-initial-iteration}
Let
\[
    F:
    \mathbf{FDom}^{\triangleleft}_{\mathbb V,\bot}
    \to
    \mathbf{FDom}^{\triangleleft}_{\mathbb V,\bot}
\]
be a pointed projection-pair endopseudofunctor.

The \emph{pointed initial \(F\)-iteration} is the sequential split
projection-pair system
\[
    D:\omega_{\mathbb V}
      \rightsquigarrow
      \mathbf{FDom}^{\triangleleft}_{\mathbb V,\bot}
\]
defined on objects by
\[
    D_0=\mathbf 1_\bot,
    \qquad
    D_{n+1}=F(D_n).
\]
The adjacent transition projection pair
\[
    D_n\triangleleft D_{n+1}
\]
is obtained by applying \(F^n\) to the canonical projection pair
\[
    \mathbf 1_\bot\triangleleft F(\mathbf 1_\bot).
\]
For \(n\leq m\), the transition projection pair
\[
    D_n\triangleleft D_m
\]
is the composite of the adjacent transition projection pairs.
\end{definition}

\begin{definition}[Pointed generated invariant]
\label{def:pointed-generated-invariant}
Let
\[
    F:
    \mathbf{FDom}^{\triangleleft}_{\mathbb V,\bot}
    \to
    \mathbf{FDom}^{\triangleleft}_{\mathbb V,\bot}
\]
be a pointed projection-pair endopseudofunctor, and let
\[
    D:\omega_{\mathbb V}
      \rightsquigarrow
      \mathbf{FDom}^{\triangleleft}_{\mathbb V,\bot}
\]
be its pointed initial iteration.

A \emph{pointed generated invariant} of \(F\) consists of a pointed fuzzy domain
\(X\), an equivalence
\[
    X\simeq F(X),
\]
and a compatible projection-pair cone
\[
    (\varepsilon_n,\pi_n):D_n\triangleleft X
    \qquad(n\in\omega_{\mathbb V})
\]
which exhibits \(X\) as the \(W_\omega\)-bilimit of the pointed initial
iteration \(D\).
\end{definition}

\begin{theorem}[Pointed finite-flat minimal-invariant theorem]
\label{thm:pointed-ff-minimal-invariant}
Assume that \(\mathbb V\) is sequentially finite-flat cofinal.  Let
\[
    F:
    \mathbf{FDom}^{\triangleleft}_{\mathbb V,\bot}
    \to
    \mathbf{FDom}^{\triangleleft}_{\mathbb V,\bot}
\]
be a pointed projection-pair endopseudofunctor preserving \(W_\omega\)-bilimits
of pointed sequential projection-pair systems.

Let
\[
    D:\omega_{\mathbb V}
      \rightsquigarrow
      \mathbf{FDom}^{\triangleleft}_{\mathbb V,\bot}
\]
be the pointed initial \(F\)-iteration.  Suppose \(D\) has a
\(W_\omega\)-bilimit
\[
    D_\infty.
\]
Then there is a canonical equivalence of pointed fuzzy domains
\[
    D_\infty\simeq F(D_\infty).
\]

Moreover, \(D_\infty\) is unique as a pointed generated invariant: if \(X\) is
any other pointed generated invariant of \(F\), then there is an equivalence
\[
    X\simeq D_\infty
\]
unique up to invertible compatible \(2\)-cell.
\end{theorem}

\begin{proof}
Let
\[
    s:\omega_{\mathbb V}\to\omega_{\mathbb V}
\]
be the successor endofunctor.  By construction of the pointed initial
iteration,
\[
    Ds\cong FD
\]
as split projection-pair systems.

Since \(s\) is \(W_\omega\)-cofinal and bilimit-cofinal for \(D\),
Theorem~\ref{thm:shifted-bilimits-ff} shows that
\[
    D_\infty
\]
is a \(W_\omega\)-bilimit of
\[
    Ds.
\]
Since \(F\) preserves \(W_\omega\)-bilimits,
\[
    F(D_\infty)
\]
is a \(W_\omega\)-bilimit of
\[
    FD.
\]
Using
\[
    Ds\cong FD,
\]
both \(D_\infty\) and \(F(D_\infty)\) are \(W_\omega\)-bilimits of the same
system.  By uniqueness,
\[
    D_\infty\simeq F(D_\infty).
\]

If \(X\) is another pointed generated invariant, then \(X\) is also a
\(W_\omega\)-bilimit of the pointed initial iteration \(D\).  Hence uniqueness
of \(W_\omega\)-bilimits gives
\[
    X\simeq D_\infty
\]
uniquely up to invertible compatible \(2\)-cell.
\end{proof}

\section{Pointed finite-flat generated algebraic compactness}
\label{sec:pointed-ff-algebraic-compactness}

The minimal-invariant theorem gives an equivalence
\[
    D_\infty\simeq F(D_\infty).
\]
To express the algebraic-compactness content, we record the associated
initial-algebra and final-coalgebra behavior inside the generated
projection-pair fragment.

Throughout this section \(F\) is assumed to act locally on fuzzy-domain
morphisms and to preserve finite-flat colimits on hom fuzzy domains.

\begin{definition}[Locally finite-flat-continuous pointed endopseudofunctor]
\label{def:locally-ff-continuous-endopseudofunctor}
A pointed projection-pair endopseudofunctor
\[
    F:
    \mathbf{FDom}^{\triangleleft}_{\mathbb V,\bot}
    \to
    \mathbf{FDom}^{\triangleleft}_{\mathbb V,\bot}
\]
is \emph{locally finite-flat-continuous} if it is the restriction to projection
pairs of a pseudofunctor
\[
    \widetilde F:
    \mathbf{FDom}^{(2)}_{\mathbb V,\bot}
    \to
    \mathbf{FDom}^{(2)}_{\mathbb V,\bot}
\]
on the locally ordinary \(2\)-category of pointed fuzzy domains and
fuzzy-domain morphisms, and for every pair of pointed fuzzy domains \(X,Y\),
the induced hom-action
\[
    \widetilde F_{X,Y}:
    [X,Y]_{\mathrm{ff}}
    \to
    [FX,FY]_{\mathrm{ff}}
\]
preserves finite-flat weighted colimits.
\end{definition}

\begin{definition}[Generated projection-pair \(F\)-algebra]
\label{def:generated-projection-pair-F-algebra}
Let
\[
    F:
    \mathbf{FDom}^{\triangleleft}_{\mathbb V,\bot}
    \to
    \mathbf{FDom}^{\triangleleft}_{\mathbb V,\bot}
\]
be locally finite-flat-continuous, and let
\[
    D:\omega_{\mathbb V}\rightsquigarrow
    \mathbf{FDom}^{\triangleleft}_{\mathbb V,\bot}
\]
be its pointed initial iteration.

A \emph{generated projection-pair \(F\)-algebra} is a pointed fuzzy domain
\(X\) equipped with a projection pair
\[
    a=(a_e,a_p):FX\triangleleft X.
\]
The embedding
\[
    a_e:FX\to X
\]
is the algebra structure, and the projection
\[
    a_p:X\to FX
\]
is the associated coalgebra structure.

The projection pair \(a\) generates a compatible cone
\[
    \lambda^a_n:D_n\triangleleft X
\]
by recursion:
\[
    \lambda^a_0:\mathbf 1_\bot\triangleleft X
\]
is the canonical projection pair, and
\[
    \lambda^a_{n+1}
    =
    a\circ F(\lambda^a_n).
\]
\end{definition}

\begin{definition}[Generated algebra and coalgebra categories]
\label{def:generated-algebra-coalgebra-categories}
The locally ordinary \(2\)-category of generated projection-pair \(F\)-algebras
is denoted
\[
    \operatorname{Alg}^{\mathrm{gen}}_{\triangleleft}(F).
\]
The same data may be read coalgebraically: the projection half
\[
    a_p:X\to FX
\]
of
\[
    a=(a_e,a_p):FX\triangleleft X
\]
is an \(F\)-coalgebra structure.  With the same objects, but with the
coalgebra-half compatibility emphasized, we write
\[
    \operatorname{Coalg}^{\mathrm{gen}}_{\triangleleft}(F).
\]
\end{definition}

\begin{construction}[The canonical generated algebra-coalgebra]
\label{con:canonical-generated-algebra-coalgebra}
Let
\[
    D_\infty
\]
be the \(W_\omega\)-bilimit of the pointed initial iteration.  By
Theorem~\ref{thm:pointed-ff-minimal-invariant}, there is an equivalence
\[
    D_\infty\simeq F(D_\infty).
\]
Choosing the corresponding projection-pair representative gives a canonical
generated projection-pair algebra
\[
    a_\infty=(a_{\infty,e},a_{\infty,p}):
    F(D_\infty)\triangleleft D_\infty.
\]
\end{construction}

\begin{lemma}[Generated folds from the bilimit]
\label{lem:generated-folds-from-bilimit}
Let
\[
    a=(a_e,a_p):FX\triangleleft X
\]
be a generated projection-pair \(F\)-algebra.  Then the generated cone
\[
    \lambda^a_n:D_n\triangleleft X
\]
induces a projection pair
\[
    \operatorname{fold}_a:D_\infty\triangleleft X.
\]
Moreover, its embedding half is an algebra morphism
\[
    (D_\infty,a_\infty)\to(X,a),
\]
and its projection half is a coalgebra morphism
\[
    (X,a_p)\to(D_\infty,a_{\infty,p}).
\]
\end{lemma}

\begin{proof}
Since \(D_\infty\) represents compatible cones over the pointed initial
iteration, the generated compatible cone
\[
    \lambda^a_n:D_n\triangleleft X
\]
is represented by a projection pair
\[
    \operatorname{fold}_a:D_\infty\triangleleft X.
\]
The algebra and coalgebra compatibility equations are checked on finite
approximants.  At stage \(n+1\), both sides are obtained from stage \(n\) by
applying \(F\) and then composing with \(a\).  Detection by the bilimit
representation gives the required compatible \(2\)-cells.
\end{proof}

\begin{lemma}[Uniqueness of generated folds]
\label{lem:uniqueness-generated-folds}
Let
\[
    a=(a_e,a_p):FX\triangleleft X
\]
be a generated projection-pair \(F\)-algebra.  Any two generated algebra
morphisms
\[
    (D_\infty,a_\infty)\to(X,a)
\]
are uniquely isomorphic by an invertible compatible \(2\)-cell.  Dually, any
two generated coalgebra morphisms
\[
    (X,a_p)\to(D_\infty,a_{\infty,p})
\]
are uniquely isomorphic by an invertible compatible \(2\)-cell.
\end{lemma}

\begin{proof}
The restrictions of a generated algebra morphism to finite approximants are
forced recursively.  At stage \(0\) they are the canonical map from
\(\mathbf 1_\bot\).  If the restrictions agree at stage \(n\), algebra
compatibility forces agreement at stage \(n+1\).  Thus any two such morphisms
represent the same compatible cone, and uniqueness follows from the
representing property of the bilimit.  The coalgebra statement is the mate
dual.
\end{proof}

\begin{theorem}[Pointed finite-flat generated algebraic compactness]
\label{thm:pointed-ff-algebraic-compactness}
Assume that \(\mathbb V\) is sequentially finite-flat cofinal.  Let
\[
    F:
    \mathbf{FDom}^{\triangleleft}_{\mathbb V,\bot}
    \to
    \mathbf{FDom}^{\triangleleft}_{\mathbb V,\bot}
\]
be a pointed projection-pair endopseudofunctor which preserves
\(W_\omega\)-bilimits of pointed sequential projection-pair systems and is
locally finite-flat-continuous.

Let
\[
    D:\omega_{\mathbb V}
      \rightsquigarrow
      \mathbf{FDom}^{\triangleleft}_{\mathbb V,\bot}
\]
be the pointed initial \(F\)-iteration, and suppose it has a
\(W_\omega\)-bilimit
\[
    D_\infty.
\]
Then the canonical projection pair
\[
    a_\infty=(a_{\infty,e},a_{\infty,p}):
    F(D_\infty)\triangleleft D_\infty
\]
exhibits \(D_\infty\) as the generated algebraically compact solution of \(F\):

\begin{enumerate}
    \item the embedding half
    \[
        a_{\infty,e}:F(D_\infty)\to D_\infty
    \]
    is initial in
    \[
        \operatorname{Alg}^{\mathrm{gen}}_{\triangleleft}(F);
    \]

    \item the projection half
    \[
        a_{\infty,p}:D_\infty\to F(D_\infty)
    \]
    is final in
    \[
        \operatorname{Coalg}^{\mathrm{gen}}_{\triangleleft}(F);
    \]

    \item the algebra and coalgebra structures are inverse halves of the same
    projection-pair equivalence
    \[
        F(D_\infty)\triangleleft D_\infty.
    \]
\end{enumerate}
\end{theorem}

\begin{proof}
Theorem~\ref{thm:pointed-ff-minimal-invariant} gives
\[
    D_\infty\simeq F(D_\infty).
\]
Choosing the corresponding projection-pair representative gives
\[
    a_\infty:F(D_\infty)\triangleleft D_\infty.
\]

For any generated projection-pair \(F\)-algebra
\[
    a:FX\triangleleft X,
\]
Lemma~\ref{lem:generated-folds-from-bilimit} gives a generated fold
\[
    \operatorname{fold}_a:D_\infty\triangleleft X.
\]
Its embedding half is an algebra morphism, and its projection half is a
coalgebra morphism.  Lemma~\ref{lem:uniqueness-generated-folds} gives
uniqueness up to invertible compatible \(2\)-cell.  Hence the embedding half of
\(a_\infty\) is initial in the generated algebra category, and the projection
half is final in the generated coalgebra category.

The two structures are inverse halves of the same projection-pair equivalence
by construction.
\end{proof}

\section{The crisp separated case}
\label{sec:crisp-separated-ff-rde}

For
\[
    \Lambda=2,
\]
finite-flat weights are ordinary ideals.  On separated fuzzy domains, strict
algebras for the ideal monad are precisely dcpos, and fuzzy-domain morphisms
are precisely Scott-continuous maps.

Let
\[
    \omega=0\leq1\leq2\leq\cdots
\]
and let
\[
    W_\omega:\omega^{\op}\to 2
\]
be the constant-true weight:
\[
    W_\omega(n)=1.
\]
This is the ordinary directed ideal \(\omega\subseteq\omega\).

For a dcpo system
\[
    D_0\triangleleft D_1\triangleleft D_2\triangleleft\cdots
\]
of projection pairs, a \(W_\omega\)-bilimit is exactly the usual bilimit:
there are projection pairs
\[
    (\varepsilon_n,\pi_n):D_n\triangleleft D_\infty
\]
such that the approximants satisfy
\[
    \bigvee_{n<\omega}
        \varepsilon_n\pi_n
    =
    \id_{D_\infty}
\]
in the dcpo of Scott-continuous endomorphisms of \(D_\infty\).

In the crisp case, the ordinary diagonal
\[
    \Delta_\omega:\omega\to\omega\times\omega,
    \qquad
    n\longmapsto(n,n),
\]
is a monotone map.  It is cofinal for the ideal weight because for every pair
\[
    (m,n)\in\omega\times\omega
\]
there is
\[
    k\geq m,n.
\]
Hence
\[
    (\Delta_\omega)_!W_\omega
    =
    W_\omega\boxtimes W_\omega.
\]
This diagonal is the cartesian feature that permits same-index function-space
equations.

\begin{theorem}[Crisp separated same-index function-space collapse]
\label{thm:crisp-separated-same-index-collapse}
Let
\[
    \Lambda=2.
\]
Let
\[
    D_0\triangleleft D_1\triangleleft D_2\triangleleft\cdots
\]
be a sequential dcpo projection-pair system with bilimit \(D_\infty\), and
suppose
\[
    D_{n+1}\cong [D_n,D_n]
\]
as projection-pair systems.  Then
\[
    D_\infty\cong [D_\infty,D_\infty].
\]
\end{theorem}

\begin{proof}
For \(\Lambda=2\), the finite-flat weight
\[
    W_\omega:\omega^{\op}\to 2
\]
is the ordinary directed ideal.  The product weight
\[
    W_\omega\boxtimes W_\omega
\]
is the ordinary directed ideal on \(\omega\times\omega\).

The diagonal
\[
    \Delta_\omega:\omega\to\omega\times\omega
\]
is cofinal, so restriction along \(\Delta_\omega\) preserves the relevant
bilimit.  The external function-space system is
\[
    (m,n)\longmapsto [D_n,D_m].
\]
Restricting along the diagonal gives
\[
    n\longmapsto [D_n,D_n]\cong D_{n+1}.
\]
The successor shift is cofinal, so the shifted chain has the same bilimit
\(D_\infty\).  On the other hand, the standard bilimit theorem for dcpo
projection-pair chains gives that
\[
    [D_\infty,D_\infty]
\]
is the bilimit of the same external function-space system.  By uniqueness of
bilimits,
\[
    D_\infty\cong [D_\infty,D_\infty].
\]
\end{proof}

Similarly, for a locally Scott-continuous bilimit-preserving endofunctor \(T\)
on dcpos, the effectful constructor
\[
    \Phi_T(D)=[D,TD]
\]
gives, after diagonal collapse, the equation
\[
    D_\infty\cong [D_\infty,TD_\infty].
\]
For a bilimit-preserving coeffect functor \(G\), the coeffectful form is
\[
    D_\infty\cong [GD_\infty,D_\infty],
\]
and for mixed effect/coeffect structure,
\[
    D_\infty\cong [GD_\infty,TD_\infty].
\]

In the crisp separated case the Kelly tensor product of fuzzy domains is the
cartesian product of dcpos, so the application maps have their usual cartesian
types:
\[
    D_\infty\times D_\infty\to D_\infty,
\]
\[
    D_\infty\times D_\infty\to TD_\infty,
\]
and
\[
    D_\infty\times GD_\infty\to TD_\infty
\]
in the pure, effectful, and mixed cases respectively.

\begin{proposition}[Crisp separated pointed minimal invariants]
\label{prop:crisp-separated-pointed-minimal-invariants}
For
\[
    \Lambda=2,
\]
the pointed finite-flat minimal-invariant and generated algebraic-compactness
theorems specialize to the classical pointed-dcpo bilimit construction.

More explicitly:

\begin{enumerate}
    \item pointed fuzzy domains are precisely pointed dcpos;

    \item fuzzy projection pairs are precisely embedding-projection pairs
    \[
        e:D\to E,
        \qquad
        p:E\to D,
        \qquad
        pe=\id_D,
        \qquad
        ep\leq\id_E;
    \]

    \item the finite-flat weight
    \[
        W_\omega:\omega^{\op}\to 2
    \]
    is the ordinary directed ideal on \(\omega\);

    \item the \(W_\omega\)-bilimit approximation identity becomes
    \[
        \bigvee_{n<\omega}\varepsilon_n\pi_n
        =
        \id_{D_\infty}
    \]
    in the dcpo of Scott-continuous endomorphisms of \(D_\infty\);

    \item a locally finite-flat-continuous endopseudofunctor becomes a locally
    Scott-continuous endofunctor on the projection-pair category of pointed
    dcpos;

    \item the canonical projection pair
    \[
        F(D_\infty)\triangleleft D_\infty
    \]
    gives the generated coincidence of the initial algebra and final coalgebra.
\end{enumerate}
\end{proposition}

\begin{proof}
For \(\Lambda=2\), finite-flat weights are ordinary ideals and strict algebras
for the ideal monad on separated preorders are precisely dcpos.  An enriched
initial object is exactly a least element, so pointed fuzzy domains are pointed
dcpos.

A fuzzy projection pair
\[
    (e,p):D\triangleleft E
\]
is an adjunction
\[
    e\dashv p
\]
with invertible unit.  In order-enriched notation this is precisely
\[
    pe=\id_D,
    \qquad
    ep\leq\id_E.
\]
Thus projection pairs are the usual embedding-projection pairs.

The weight
\[
    W_\omega:\omega^{\op}\to 2
\]
is constant true, hence is the ordinary directed ideal on \(\omega\).  Therefore
a \(W_\omega\)-weighted colimit in an internal hom dcpo is an ordinary directed
join.  The approximation identity
\[
    W_\omega*(n\mapsto\varepsilon_n\pi_n)
    \cong
    \id_{D_\infty}
\]
therefore becomes
\[
    \bigvee_{n<\omega}\varepsilon_n\pi_n
    =
    \id_{D_\infty}.
\]

The generated projection-pair algebra
\[
    F(D_\infty)\triangleleft D_\infty
\]
has embedding half
\[
    F(D_\infty)\to D_\infty
\]
and projection half
\[
    D_\infty\to F(D_\infty).
\]
These are the classical inverse algebra and coalgebra maps obtained from the
bilimit of the initial embedding-projection chain.
\end{proof}

\section{Summary}
\label{sec:summary-ff-rde}

Recursive domain equations in fuzzy domains are governed by the finite-flat
doctrine.  The primitive approximation objects are finite-flat approximation
shapes
\[
    (A,W),
    \qquad
    W\in\operatorname{Idl}_{\mathbb V}(A).
\]
Cofinality of weights is expressed by lower direct image:
\[
    h:(A,W)\to(B,U)
    \quad\Longleftrightarrow\quad
    h_!W\cong U.
\]
An iteration shape is a finite-flat approximation shape equipped with a
weight-cofinal self-map
\[
    s:A\to A,
    \qquad
    s_!W\cong W.
\]

Tails are canonical support-comma categories
\[
    \operatorname{Tail}_A(a_0),
\]
equipped with their canonical projection
\[
    \tau_{a_0}:\operatorname{Tail}_A(a_0)\to A.
\]
Tail-cofinality says that
\[
    \tau_{a_0}^*W
\]
is finite-flat and
\[
    (\tau_{a_0})_!\tau_{a_0}^{*}W\cong W.
\]
Common successors are encoded by
\[
    \operatorname{Com}(P),
\]
whose objects are tuples
\[
    (i,j,k,u:i\to k,v:j\to k).
\]
The canonical maps
\[
    c_P:\operatorname{Com}(P)\to P\tensor P,
    \qquad
    m_P:\operatorname{Com}(P)\to P
\]
are cofinal when
\[
    (c_P)_!m_P^*R\cong R\boxtimes R,
    \qquad
    (m_P)_!m_P^*R\cong R.
\]

For bilimits, weight-cofinality is paired with bilimit-cofinality: restriction
along the indexing map must preserve compatible projection-pair cone
categories.  With this distinction in place, a split projection-pair system
\[
    D:A\rightsquigarrow\mathbf{FDom}_{\mathbb V}^{\triangleleft}
\]
satisfying
\[
    Ds\cong FD
\]
has a generated solution whenever \(D\) has a \(W\)-bilimit, \(s\) is
bilimit-cofinal for \(D\), and \(F\) preserves that bilimit:
\[
    D_\infty\simeq F(D_\infty).
\]

The general internal-hom theorem is external-product-indexed.  If
\[
    X_\infty
    \qquad\text{and}\qquad
    Y_\infty
\]
are \(W\)-bilimits of split systems \(X\) and \(Y\), then
\[
    [X_\infty,Y_\infty]_{\mathrm{ff}}
\]
has external-product approximants
\[
    (a,a')\longmapsto [X_{a'},Y_a]_{\mathrm{ff}},
\]
and these satisfy
\[
    (W\boxtimes W)*
    \bigl((a,a')\mapsto E_{a,a'}\Pi_{a,a'}\bigr)
    \cong
    \id_{[X_\infty,Y_\infty]_{\mathrm{ff}}}.
\]
Compatible cones over the external-product system reconstruct canonically to
lax application data.  They do not generally reconstruct to projection pairs,
because the common-successor comparison supplies
\[
    \rho\delta\Rightarrow\id
\]
rather than an invertible unit
\[
    \id\Rightarrow\rho\delta.
\]
This is exactly the resource-sensitive obstruction caused by the absence of a
canonical diagonal
\[
    A\to A\tensor A.
\]

In the pointed case, a sequential finite-flat approximation shape
\[
    (\omega_{\mathbb V},W_\omega,s)
\]
replaces the ordinary chain
\[
    0\leq1\leq2\leq\cdots.
\]
The pointed initial iteration
\[
    D_0=\mathbf 1_\bot,
    \qquad
    D_{n+1}=F(D_n)
\]
has a \(W_\omega\)-bilimit \(D_\infty\).  If \(F\) preserves this bilimit, then
\[
    D_\infty\simeq F(D_\infty).
\]
If \(F\) is locally finite-flat-continuous, the resulting projection pair
\[
    F(D_\infty)\triangleleft D_\infty
\]
is the generated algebraically compact solution: its embedding half is the
initial generated algebra, its projection half is the final generated
coalgebra, and the two structures are inverse halves of the same projection-pair
equivalence.

In the crisp separated case, taking the ordinary ideal
\[
    W_\omega=\omega\subseteq\omega
\]
recovers the standard bilimit method for solving recursive domain equations of
dcpos.  The ordinary reflexive-domain equation
\[
    D_\infty\cong [D_\infty,D_\infty]
\]
is recovered precisely because the cartesian diagonal
\[
    \omega\to\omega\times\omega
\]
collapses the external-product function-space approximation to a same-index
chain.